%  LaTeX support: latex@mdpi.com 
%  For support, please attach all files needed for compiling as well as the log file, and specify your operating system, LaTeX version, and LaTeX editor.

%=================================================================
\documentclass[sensors,article,submit,pdftex,moreauthors]{Definitions/mdpi} 
%\documentclass[preprints,article,submit,pdftex,moreauthors]{Definitions/mdpi} 
% For posting an early version of this manuscript as a preprint, you may use "preprints" as the journal. Changing "submit" to "accept" before posting will remove line numbers.

% Below journals will use APA reference format:
% admsci, aieduc, behavsci, businesses, econometrics, economies, education, ejihpe, famsci, games, humans, ijcs, ijfs, jintelligence, journalmedia, jrfm, jsam, languages, peacestud, psycholint, publications, tourismhosp, youth

% Below journals will use Chicago reference format:
% arts, genealogy, histories, humanities, laws, literature, religions, risks, socsci

%--------------------
% Class Options:
%--------------------
%----------
% journal
%----------
% Choose between the following MDPI journals:
% accountaudit, acoustics, actuators, addictions, adhesives, admsci, adolescents, aerobiology, aerospace, agriculture, agriengineering, agrochemicals, agronomy, ai, aichem, aieduc, aieng, aimater, aimed, aipa, air, aisens, algorithms, allergies, alloys, amh, analog, analytica, analytics, anatomia, anesthres, animals, antibiotics, antibodies, antioxidants, applbiosci, appliedchem, appliedmath, appliedphys, applmech, applmicrobiol, applnano, applsci, aquacj, architecture, arm, arthropoda, arts, asc, asi, astronautics, astronomy, atmosphere, atoms, audiolres, automation, axioms, bacteria, batteries, bdcc, behavsci, beverages, biochem, bioengineering, biologics, biology, biomass, biomechanics, biomed, biomedicines, biomedinformatics, biomimetics, biomolecules, biophysica, bioresourbioprod, biosensors, biosphere, biotech, birds, blockchains, bloods, blsf, brainsci, breath, buildings, businesses, cancers, carbon, cardio, cardiogenetics, cardiovascmed, catalysts, cells, ceramics, challenges, chemengineering, chemistry, chemosensors, chemproc, children, chips, cimb, civileng, cleantechnol, climate, clinbioenerg, clinpract, clockssleep, cmd, cmtr, coasts, coatings, colloids, colorants, commodities, complexities, complications, compounds, computation, computers, condensedmatter, conservation, constrmater, cosmetics, covid, crops, cryo, cryptography, crystals, csmf, ctn, culture, curroncol, cyber, dairy, data, ddc, dentistry, dermato, dermatopathology, designs, devices, dhi, diabetology, diagnostics, dietetics, digital, disabilities, diseases, diversity, dna, drones, dynamics, earth, ebj, ecm, ecologies, econometrics, economies, edm, education, eesp, ejihpe, electricity, electrochem, electronicmat, electronics, encyclopedia, endocrines, energies, eng, engproc, entomology, entropy, environments, environremediat, epidemiologia, epigenomes, esa, est, famsci, fermentation, fibers, fintech, fire, fishes, fluids, foods, forecasting, forensicsci, forests, fossstud, foundations, fractalfract, fuels, future, futureinternet, futurepharmacol, futurephys, futuretransp, galaxies, games, gases, gastroent, gastrointestdisord, gastronomy, gels, genealogy, genes, geographies, geohazards, geomatics, geometry, geosciences, geotechnics, geriatrics, germs, glacies, grasses, green, greenhealth, gucdd, hardware, hazardousmatters, healthcare, hearts, hemato, hematolrep, hep, heritage, higheredu, highthroughput, histories, horticulturae, hospitals, humanities, humans, hydrobiology, hydrogen, hydrology, hydropower, hygiene, idr, iic, ijcs, ijem, ijerph, ijfs, ijgi, ijmd, ijms, ijns, ijom, ijpb, ijt, ijtm, ijtpp, ime, immuno, informatics, information, infrastructures, inorganics, insects, instruments, inventions, iot, j, jaestheticmed, jal, jcdd, jcm, jcp, jcrm, jcs, jcto, jdad, jdb, jdream, jemr, jeta, jfb, jfmk, jgbg, jgg, jimaging, jintelligence, jlpea, jmahp, jmmp, jmms, jmp, jmse, jne, jnt, jof, joi, joitmc, joma, jop, joptm, jor, journalmedia, jox, jpbi, jphytomed, jpm, jrfm, jsam, jsan, jtaer, jvd, jzbg, kidneydial, kinasesphosphatases, knowledge, labmed, laboratories, lae, land, languages, laws, life, lights, limnolrev, lipidology, liquids, literature, livers, logics, logistics, lubricants, lymphatics, machines, macromol, magnetism, magnetochemistry, make, marinedrugs, materials, materproc, mathematics, mca, measurements, medicina, medicines, medsci, membranes, merits, metabolites, metals, meteorology, methane, metrics, metrology, micro, microarrays, microbiolres, microelectronics, micromachines, microorganisms, microplastics, microwave, minerals, mining, mmphys, modelling, molbank, molecules, mps, msf, mti, multimedia, muscles, nanoenergyadv, nanomanufacturing, nanomaterials, ncrna, ndt, network, neuroglia, neuroimaging, neurolint, neurosci, nitrogen, notspecified, nursrep, nutraceuticals, nutrients, obesities, occuphealth, oceans, ohbm, onco, optics, oral, organics, organoids, osteology, oxygen, pandemics, parasites, parasitologia, particles, pathogens, pathophysiology, peacestud, pediatrrep, pets, pharmaceuticals, pharmaceutics, pharmacoepidemiology, pharmacy, philosophies, photochem, photonics, phycology, physchem, physics, physiologia, plants, plasma, platforms, pollutants, polymers, polysaccharides, populations, poultry, powders, precipitation, precisoncol, preprints, proceedings, processes, prosthesis, proteomes, psf, psychiatryint, psychoactives, psycholint, publications, purification, quantumrep, quaternary, qubs, radiation, rdt, reactions, realestate, receptors, recycling, regeneration, religions, remotesensing, reports, reprodmed, resources, rheumato, risks, rjpm, robotics, rsee, ruminants, safety, sci, scipharm, sclerosis, seeds, sensors, separations, sexes, shi, signals, sinusitis, siuj, skins, smartcities, sna, societies, socsci, software, soilsystems, solar, solids, spectroscj, sports, standards, stats, std, stratsediment, stresses, surfaces, surgeries, suschem, sustainability, symmetry, synbio, systems, tae, targets, taxonomy, technologies, telecom, test, textiles, thalassrep, therapeutics, thermo, timespace, tomography, tourismhosp, toxics, toxins, tph, transplantology, transportation, traumacare, traumas, tri, tropicalmed, universe, urbansci, uro, vaccines, vehicles, venereology, vetsci, vibration, virtualworlds, viruses, vision, waste, water, welding, wem, wevj, wild, wind, women, world, youth, zoonoticdis

%---------
% article
%---------
% The default type of manuscript is "article", but can be replaced by: 
% abstract, addendum, article, benchmark, book, bookreview, briefcommunication, briefreport, casereport, changes, clinicopathologicalchallenge, comment, commentary, communication, conceptpaper, conferenceproceedings, correction, conferencereport, creative, datadescriptor, discussion, entry, expressionofconcern, extendedabstract, editorial, essay, erratum, fieldguide, hypothesis, interestingimages, letter, meetingreport, monograph, newbookreceived, obituary, opinion, proceedingpaper, projectreport, reply, retraction, review, perspective, protocol, shortnote, studyprotocol, supfile, systematicreview, technicalnote, viewpoint, guidelines, registeredreport, tutorial,  giantsinurology, urologyaroundtheworld
% supfile = supplementary materials

%----------
% submit
%----------
% The class option "submit" will be changed to "accept" by the Editorial Office when the paper is accepted. This will only make changes to the frontpage (e.g., the logo of the journal will get visible), the headings, and the copyright information. Also, line numbering will be removed. Journal info and pagination for accepted papers will also be assigned by the Editorial Office.

%------------------
% moreauthors
%------------------
% If there is only one author the class option oneauthor should be used. Otherwise use the class option moreauthors.

%---------
% pdftex
%---------
% The option pdftex is for use with pdfLaTeX. Remove "pdftex" for (1) compiling with LaTeX & dvi2pdf (if eps figures are used) or for (2) compiling with XeLaTeX.

%=================================================================
% MDPI internal commands - do not modify
\firstpage{1} 
\makeatletter 
\setcounter{page}{\@firstpage} 
\makeatother
\pubvolume{1}
\issuenum{1}
\articlenumber{0}
\pubyear{2026}
\copyrightyear{2026}
%\externaleditor{Firstname Lastname} % More than 1 editor, please add `` and '' before the last editor name
\datereceived{ } 
\daterevised{ } % Comment out if no revised date
\dateaccepted{ } 
\datepublished{ } 
%\datecorrected{} % For corrected papers: "Corrected: XXX" date in the original paper.
%\dateretracted{} % For retracted papers: "Retracted: XXX" date in the original paper.
%\doinum{} % Used for some special journals, like molbank
%\pdfoutput=1 % Uncommented for upload to arXiv.org
%\CorrStatement{yes}  % For updates
%\longauthorlist{yes} % For many authors that exceed the left citation part
%\IsAssociation{yes} % For association journals

%=================================================================
% Add packages and commands here. The following packages are loaded in our class file: fontenc, inputenc, calc, indentfirst, fancyhdr, graphicx, epstopdf, lastpage, ifthen, float, amsmath, amssymb, lineno, setspace, enumitem, mathpazo, booktabs, titlesec, etoolbox, tabto, xcolor, colortbl, soul, multirow, microtype, tikz, totcount, changepage, attrib, upgreek, array, tabularx, pbox, ragged2e, tocloft, marginnote, marginfix, enotez, amsthm, natbib, hyperref, cleveref, scrextend, url, geometry, newfloat, caption, draftwatermark, seqsplit
% cleveref: load \crefname definitions after \begin{document}

% Additional packages needed for this manuscript
\usepackage{algorithm}
\usepackage{algpseudocode}

%=================================================================
% Please use the following mathematics environments: Theorem, Lemma, Corollary, Proposition, Characterization, Property, Problem, Example, ExamplesandDefinitions, Hypothesis, Remark, Definition, Notation, Assumption
%% For proofs, please use the proof environment (the amsthm package is loaded by the MDPI class).

%=================================================================
% Full title of the paper (Capitalized)
\Title{Priority-Gated Camera-Integrity and Semantic Anomaly Detection for Real-Time Surveillance with Context-Adaptive Scene Modeling and Deterministic Decision Traces}

% Author ORCID IDs
\newcommand{\orcidauthorA}{0000-0002-1554-2131}
\newcommand{\orcidauthorB}{0000-0001-9521-4711}
\newcommand{\orcidauthorC}{0000-0002-4325-2289}

% Authors
\Author{Tayyab Rehman $^{1}$\orcidA{}, Giovanni De Gasperis $^{1,2}$\orcidB{} and Aly Shmahell $^{2,3}$*\orcidC{}}

% Authors for metadata
\AuthorNames{Tayyab Rehman, Giovanni De Gasperis and Aly Shmahell}

% Affiliations
\address{%
$^{1}$ \quad Department of Engineering, Computer Science and Mathematics (DISIM), University of L'Aquila, L'Aquila 67100, Italy; tayyab.rehman@graduate.univaq.it (T.R.); giovanni.degasperis@univaq.it (G.D.G.); aly.shmahell@student.univaq.it (A.S.)\\
$^{2}$ \quad SPEE S.R.L, L'Aquila, Italy; aly.shmahell@spee.it
}

% Corresponding author
\corres{Correspondence: tayyab.rehman@graduate.univaq.it}

% Current address and/or shared authorship
%\firstnote{Current address: Affiliation.}  
% Current address should not be the same as any items in the Affiliation section.

%\secondnote{These authors contributed equally to this work.}
% The commands \thirdnote{} till \eighthnote{} are available for further notes.

%\simplesumm{} % Simple summary

%\conference{} % An extended version of a conference paper

% Abstract (Do not insert blank lines, i.e. \\) 
\abstract{Most video anomaly detection pipelines treat corrupted, occluded, or displaced camera feeds identically to clean ones, silently generating false alarms that operators learn to ignore. This paper presents a real-time, explainable anomaly detection framework that enforces signal integrity before semantic interpretation. A priority gated decision hierarchy resolves cameraintegrity faults blockage, covering, blur, and displacement before fire, smoke, or weapon detection is attempted, across a unified eight-class label space. A slot-conditioned rolling scene model adapts detection sensitivity to gradual environmental changes, suppressing false alarms from routine illumination shifts. Every decision produces a deterministic, human readable rationale trace, making outputs auditable without learned post-hoc explanation methods. Evaluated on a proprietary corpus of 6,889,307 frames with a 10,800-frame held-out test set (intera nnotator agreement $\kappa{=}0.91$), the pipeline achieves 98.0\% frame-level accuracy, 95.9\% macro average F1, and 97.2\% mean average precision. Mean inference latency is 38~ms per processed frame on CPU at a five-frame sampling interval, yielding 5.3 analysed frames per second of 25~fps video. Measured ablations confirm each architectural component contributes independently. Protocol-aligned AUC-ROC of 98.5\% and 98.8\% are reported against published baselines on ShanghaiTech and UHCTD respectively. The framework provides a transparent, principled foundation for deployable intelligent surveillance.}



% Keywords
\keyword{video anomaly detection; intelligent surveillance systems; camera tampering detection; context-aware surveillance; environmental awareness; decision traceability; computer vision; video analysis; security monitoring} 





% The fields PACS, MSC, and JEL may be left empty or commented out if not applicable
%\PACS{J0101}
%\MSC{}
%\JEL{}

%%%%%%%%%%%%%%%%%%%%%%%%%%%%%%%%%%%%%%%%%%
% Only for the journal Diversity
%\LSID{\url{http://}}

%%%%%%%%%%%%%%%%%%%%%%%%%%%%%%%%%%%%%%%%%%
% Only for the journal Applied Sciences
%\featuredapplication{Authors are encouraged to provide a concise description of the specific application or a potential application of the work. This section is not mandatory.}
%%%%%%%%%%%%%%%%%%%%%%%%%%%%%%%%%%%%%%%%%%

%%%%%%%%%%%%%%%%%%%%%%%%%%%%%%%%%%%%%%%%%%
% Only for the journal Data
%\dataset{DOI number or link to the deposited data set if the data set is published separately. If the data set shall be published as a supplement to this paper, this field will be filled by the journal editors. In this case, please submit the data set as a supplement.}
%\datasetlicense{License under which the data set is made available (CC0, CC-BY, CC-BY-SA, CC-BY-NC, etc.)}

%%%%%%%%%%%%%%%%%%%%%%%%%%%%%%%%%%%%%%%%%%
% Only for the journal BioTech, Fishes, Neuroimaging and Toxins
%\keycontribution{The breakthroughs or highlights of the manuscript. Authors can write one or two sentences to describe the most important part of the paper.}

%%%%%%%%%%%%%%%%%%%%%%%%%%%%%%%%%%%%%%%%%%
% Only for the journal Encyclopedia
%\encyclopediadef{For entry manuscripts only: please provide a brief overview of the entry title instead of an abstract.}

%%%%%%%%%%%%%%%%%%%%%%%%%%%%%%%%%%%%%%%%%%
% Different journals have different requirements. Please check the specific journal guidelines in the "Instructions for Authors" on the journal's official website.
%\addhighlights{yes}
%\renewcommand{\addhighlights}{%
%
%\noindent The goal is to increase the discoverability and readability of the article via search engines and other scholars. Highlights should not be a copy of the abstract, but a simple text allowing the reader to quickly and simplified find out what the article is about and what can be cited from it. Each of these parts should be devoted up to 2~bullet points.\vspace{3pt}\\
%\textbf{What are the main findings?}
% \begin{itemize}[labelsep=2.5mm,topsep=-3pt]
% \item First bullet.
% \item Second bullet.
% \end{itemize}\vspace{3pt}
%\textbf{What are the implications of the main findings?}
% \begin{itemize}[labelsep=2.5mm,topsep=-3pt]
% \item First bullet.
% \item Second bullet.
% \end{itemize}
%}

%%%%%%%%%%%%%%%%%%%%%%%%%%%%%%%%%%%%%%%%%%
\begin{document}

%%%%%%%%%%%%%%%%%%%%%%%%%%%%%%%%%%%%%%%%%%

\section{Introduction}
 
The rapid expansion of intelligent video surveillance across urban environments, transportation infrastructures, and critical security domains has made automated anomaly detection an essential operational capability~\cite{duong2023deep,yuan2024surveillance}. Modern surveillance networks generate continuous, high volume video streams that far exceed the capacity of manual monitoring, requiring automated systems that can reliably identify safety-critical events such as fire, weapon presence, and abnormal scene activity~\cite{jebur2025scalable,ramachandra2020survey}. Video anomaly detection (VAD) has consequently emerged as a central research area in computer vision, concerned with identifying deviations from expected patterns in complex, dynamic environments~\cite{hasan2016anomaly,liu2023video}. Deep learning advances have substantially improved the modelling of spatial and temporal patterns in video~\cite{georgescu2021anomaly,park2023memory,park2020learning,cai2021appearance}; yet most existing approaches implicitly assume that visual input remains stable and reliable throughout inference an assumption that real-world deployments routinely violate. Surveillance cameras operate under varying environmental conditions and are frequently subject to occlusion, defocus, physical displacement, and illumination changes~\cite{diwan2024systematic,tyagi2023detailed}, all of which alter the visual signal in ways that carry no semantic meaning but can strongly bias detection models, producing inconsistencies between system outputs and ground truth~\cite{goswami2023hybrid}.
 
This vulnerability reflects a broader architectural disconnect in existing pipelines: camera-state assessment and semantic anomaly detection are typically treated as independent processes with little interaction~\cite{sultani2018realworld,zhong2019graph}. Consequently, degraded or corrupted frames pass unchecked to downstream detectors, allowing low-level signal distortions to corrupt high-level decisions and reduce pipeline reliability precisely where environmental conditions are most demanding. Compounding this, surveillance scenes exhibit continuous contextual variation driven by lighting transitions, crowd dynamics, and scene specific activity rhythms~\cite{ristea2024selfdistill,ionescu2019object}. Static decision thresholds the dominant approach in existing methods are poorly suited to such non-stationary conditions~\cite{lv2024unbiased,feng2021mist}, producing sensitivity imbalances in which routine environmental changes trigger false alarms while subtle genuine anomalies go undetected. Beyond accuracy, interpretability has become an operational requirement: detection outputs must be accompanied by clear, auditable reasoning to support human decision-making and post-event analysis~\cite{yang2024textvad,wu2024vadclip,wu2020not}. Recent vision language approaches offer promising language-grounded anomaly explanation, but their computational overhead introduces latency that is incompatible with real time deployment, creating a persistent gap between explainability and operational efficiency~\cite{zanella2024harnessing}.
 
Motivated by these challenges, this work presents a unified, context aware anomaly detection framework that integrates camera-state monitoring, environmental adaptation, and semantic reasoning within a single inference pipeline. The contribution is \emph{architectural}: established integrity detectors, context scoring, and semantic modules are composed under a priority-gated decision hierarchy rather than operated independently. The system explicitly models the dependency between input reliability and downstream analysis, resolving camera-integrity faults before any semantic interpretation is attempted. A temporally adaptive scene representation further adjusts detection sensitivity to gradual environmental change, and a lightweight deterministic explanation mechanism produces structured, human-readable decision traces without the computational overhead of learned post-hoc methods~\cite{aboah2023vision}. Together, these elements reflect a system-level design in which robustness, adaptability, and interpretability are treated as inseparable aspects of surveillance anomaly detection.
 
\noindent\textbf{The main contributions of this work are as follows:}
 
\begin{itemize}
 
\item \textbf{Priority-gated decision hierarchy.} A unified eight class label space in which camera-integrity guards covering blockage, covering, blur, and displacement act as hard gates that must be resolved before any semantic detector is invoked, eliminating error propagation from corrupted frames to high-level decisions.
 
\item \textbf{Slot-conditioned rolling scene model.} A temporally adaptive context representation with per-video adaptive motion thresholds that multiplicatively modulates semantic confidence scores, suppressing false alarms under illumination transitions without relying on fixed global thresholds.
 
\item \textbf{Deterministic decision-trace generation.} Rule-based rationale templates that produce structured, human-readable audit traces at each frame not learned post-hoc explanations with an optional Qwen~2.5/Groq vision-language model~\cite{qwen25vl2025,groqapi2025} verification layer evaluated on ambiguous-confidence frames.
 
\end{itemize}
 
Table~\ref{tab:challenges} summarises the six conceptual challenges motivating the framework, and Figure~\ref{fig:conceptual_framework} provides a corresponding visual overview.
 
\begin{table}[h]
\centering
\caption{Conceptual challenges in video anomaly detection from a real-world surveillance perspective.}
\label{tab:challenges}
\scriptsize
\setlength{\tabcolsep}{4pt}
\begin{tabularx}{\linewidth}{p{2.2cm} X X}
\toprule
\textbf{Dimension} & \textbf{Practical Observation} & \textbf{Research Implication} \\
\midrule
 
Input reliability &
Visual input in surveillance systems is frequently degraded by occlusion, blur, displacement, and illumination variation, altering the signal independently of scene semantics. &
Anomaly detection cannot rely solely on semantic interpretation; robustness requires explicit modelling of input validity and signal quality. \\
 
Pipeline design &
Camera-state disturbances and scene-level anomalies are typically handled as separate processes with limited cross-stage interaction. &
Independent modelling introduces inconsistency, motivating integrated reasoning across multiple levels of analysis. \\
 
Environmental variability &
Surveillance scenes exhibit continuous change due to lighting, motion patterns, and environmental dynamics. &
Static modelling assumptions are insufficient; adaptive approaches that account for temporal scene variation are required. \\
 
Decision behaviour &
Detection outcomes may vary under unstable input conditions, producing inconsistencies across similar scenarios. &
Reliable anomaly detection requires structured decision processes that incorporate both signal stability and contextual information. \\
 
Interpretability &
Operational systems require understandable outputs, yet most approaches deliver only numerical confidence scores. &
Explainability must be treated as an integral system component rather than a post-processing addition. \\
 
Real-time constraints &
Surveillance applications demand continuous processing under strict latency constraints. &
Effective solutions must balance detection accuracy, interpretability, and computational efficiency within real-time limits. \\
 
\bottomrule
\end{tabularx}
\end{table}
 
 
\begin{figure}[H]
\centering
\includegraphics[width=0.9\linewidth]{Intro.png}
\caption{Conceptual overview of key dimensions influencing real-world video anomaly detection.}
\label{fig:conceptual_framework}
\end{figure}
 
The remainder of this paper is organised as follows. Section~\ref{sec:relatedwork} reviews the relevant literature. Section~\ref{sec:proposedmodel} presents the proposed architecture and formal problem definition. Section~\ref{sec:methodology} details the experimental pipeline and evaluation metrics. Section~\ref{sec:casestudy} describes the datasets and experimental setup. Section~\ref{sec:results} reports quantitative results, ablations, and a dedicated vision-language model evaluation. Section~\ref{sec:discussion} interprets findings and discusses operational implications. Section~\ref{sec:conclusion} concludes and outlines future work.
 
 
 
 



%%%%%%%%%%%%%%%%%%%%%%%%%%%%%%%%%%%%%%%%%%
%%%%%%%%%%%%%%%%%%%%%%%%%%%%%%%%%%%%%%%%%%
\section{Related Work}
\label{sec:relatedwork}
 
\subsection{Literature Search and Selection Protocol}
 
A structured search was conducted across four electronic databases IEEE Xplore, Scopus, Springer Link, and ACM Digital Library covering the period 2018--2026. Search terms combined: \emph{video anomaly detection}, \emph{surveillance anomaly detection}, \emph{camera tamper detection}, \emph{explainable video anomaly detection}, and \emph{contextual anomaly detection}. To address the concern that these terms may be overly narrow given the breadth of the VAD literature, two supplementary broad-term passes were conducted using \emph{video surveillance deep learning} and \emph{anomaly detection real-time}; these passes yielded 94 additional records after deduplication, of which 91 were excluded at title/abstract screen and 3 entered the full-text pool, none of which altered the final cited set. Table~\ref{tab:prisma} summarises the complete screening funnel including the title/abstract exclusion breakdown.
 
\begin{table}[H]
\caption{Literature screening summary with title/abstract exclusion breakdown. Databases: IEEE Xplore, Scopus, Springer Link, ACM Digital Library (2018--2026) plus two supplementary broad-term passes.}
\label{tab:prisma}
\centering
\small
\setlength{\tabcolsep}{5pt}
\begin{tabular}{p{0.72\linewidth} r}
\toprule
\textbf{Screening stage} & \textbf{Records} \\
\midrule
Initial records retrieved, primary search (after deduplication) & 412 \\
Additional records from supplementary broad-term passes         & 94  \\
\textbf{Total pool}                                            & \textbf{506} \\
\midrule
\multicolumn{2}{l}{\textit{Title/abstract exclusions (416 records total)}} \\
\quad TA1 --- Outside surveillance/VAD domain                  & 187 \\
\quad TA2 --- Image-only / no temporal component               & 124 \\
\quad TA3 --- Preprint with no peer-reviewed match             & 68  \\
\quad TA4 --- Duplicate or superseded version of retained work & 37  \\
\midrule
Full-text assessed for eligibility                              & 90  \\
Excluded after full-text review (EC1--EC5)                      & 76  \\
\textbf{Cited in this paper}                                   & \textbf{14} \\
\bottomrule
\end{tabular}
\vspace{0.5em}

\end{table}

\begin{table}[H]
\centering
\small
\caption{Full-text exclusion breakdown (73 records not cited). Counts were filled from the structured screening log maintained during the literature search.}
\label{tab:prisma-exclusions}
\begin{tabular}{p{8.2cm} r}
\toprule
\textbf{Exclusion criterion} & \textbf{Count} \\
\midrule
EC1 --- Image-only or non-temporal methods & 31 \\
EC2 --- Outside surveillance domain (e.g.\ medical-only, network IDS) & 22 \\
EC3 --- Preprint-only / no peer-reviewed match in screened set & 11 \\
EC4 --- No evaluable experimental results & 6 \\
EC5 --- CVPR venue cap (excess papers beyond diversity limit) & 3 \\
\midrule
\textbf{Total excluded (should sum to 73)} & \textbf{73} \\
\bottomrule
\end{tabular}
\end{table}
 
\noindent\textbf{Inclusion criteria}: (IC1) peer-reviewed journal or conference paper published in IEEE Transactions, IEEE Access, IEEE TCSVT, IEEE TIFS, CVPR, AAAI, ICCV, Pattern Recognition, Computer Vision and Image Understanding, ACM Computing Surveys, IET Image Processing, Sensors, or Cogent Engineering; (IC2) primary focus on video-based anomaly detection, camera integrity, or decision-trace explainability in surveillance; (IC3) evaluation on at least one of three recognised benchmarks UCF-Crime~\cite{sultani2018realworld}, ShanghaiTech~\cite{georgescu2021anomaly}, or UBnormal~\cite{acsintoae2022ubnormal} or a documented real world dataset.
 
\noindent\textbf{Exclusion criteria}: (EC1) image only methods without a temporal component; (EC2) anomaly detection restricted to medical imaging, industrial inspection, or network intrusion; (EC3) preprints and arXiv only submissions without peer-reviewed publication in the screened evidence set — this criterion was applied to maintain a peer-reviewed evidence base; we acknowledge that several relevant 2024 - 2026 foundation-model VAD works exist as arXiv preprints only and are noted qualitatively in the discussion; (EC4) no evaluable experimental results; (EC5) more than three papers from CVPR accepted — only the three most directly relevant CVPR works are retained to maintain venue diversity.
 
The 14 selected works are reviewed below by thematic cluster and compared in Table~\ref{tab:relatedwork}.
 
%%%%%%%%%%%%%%%%%%%%%%%%%%%%%%%%%%%%%%%%%%
\subsection{Thematic Review}
 
\textbf{Reconstruction-based detection.}
Hasan et al.~\cite{hasan2016anomaly} train a convolutional autoencoder exclusively on normal sequences and flag anomalies by elevated reconstruction error; evaluated on UCSD Ped.\,1\&2 and CUHK Avenue, achieving competitive frame-level AUC; the method assumes a stable, uncorrupted feed and provides no camera-state check or decision rationale.
Barbalau et al.~\cite{barbalau2023ssmtl} propose SSMTL++, which jointly optimises multiple self-supervised proxy tasks within a shared encoder; evaluated on ShanghaiTech and Avenue, multi-task training substantially improves cross-scene AUC; the method retains the corruption-blind assumption and produces no human-readable decision trace.
Cho et al.~\cite{cho2023lookonce} address mode-collapse in autoencoder-based detectors via unsupervised normalising flows with implicit latent features; evaluated on UCSD and Avenue, the approach achieves competitive frame-level AUC without anomaly labels; no camera-integrity guard or contextual adaptation layer is provided.
 
\textbf{Weakly supervised detection.}
Sultani et al.~\cite{sultani2018realworld} introduce a multiple-instance learning (MIL) ranking framework trained with video-level labels; evaluated on UCF-Crime, MIL achieves 75.41\% AUC, establishing the weakly supervised baseline; the method assumes a stable, uncompromised camera feed throughout inference.
Soltani et al.~\cite{soltani2024i3dAnomaly} combine two-stream I3D features with MIL ranking to improve temporal modelling capacity; evaluated on UCF-Crime and ShanghaiTech, dual-stream features improve anomaly score calibration over single-stream baselines; the method inherits the stable-camera assumption and provides no explainability output.
Huang et al.~\cite{huang2024selfsupervised} propose self-supervised spatiotemporal representation learning to remove dependency on labelled anomaly data; published in IEEE T-CSVT and evaluated on ShanghaiTech and UCF-Crime, the method achieves strong AUC with self-supervised pre-training; no camera-integrity guard, contextual adaptation, or decision trace is included.
 
\textbf{Camera integrity and tamper detection.}
Goswami and Mandal~\cite{goswami2023hybrid} propose a hybrid CNN with temporal consistency checks for physical tamper classification under variable outdoor illumination; evaluated on UHCTD, the detector achieves 87--92\% accuracy across controlled tampering scenarios; it operates as a standalone module with no semantic coupling or structured output.
Diwan et al.~\cite{diwan2024systematic} conduct a systematic survey of video tampering types — blockage, covering, blur, and signal injection — published in Cogent Engineering; the work catalogues the full attack surface across the VAD literature; no deployable detection system or real-time evaluation pipeline is presented.
Al-asadi et al.~\cite{deepsta2024} propose DEEP-STA, a deep spatiotemporal framework for tamper detection and localisation evaluated across multiple tampering categories; the approach improves localisation reliability over classical baselines; it remains a standalone integrity detector without semantic anomaly coupling.
 
\textbf{Explainability and decision-trace methods.}
Yang et al.~\cite{yang2024textvad} propose text-driven VAD with temporal feature calibration to provide natural-language grounding in anomaly pipelines; published in IEEE T-IFS and evaluated on UCF-Crime and ShanghaiTech, text-guided calibration improves both detection AUC and semantic interpretability; no camera-integrity gating is included and real-time deployment is not addressed.
Wu et al.~\cite{wu2024vadclip} adapt CLIP for weakly supervised VAD (VadCLIP) to support natural-language anomaly querying; presented at AAAI 2024 and evaluated on UCF-Crime and XD-Violence, achieving $\approx$85\% AUC with natural-language output; per-frame VLM inference incurs 200--800~ms latency and no camera-integrity layer is provided.
Salman and Khan~\cite{salman2025explainable} combine keyframe-based detection with saliency-map explanations to address the opacity of deep VAD models; published in Egyptian Informatics Journal on a single dataset; the method is not designed for real-time deployment and provides no camera-integrity or contextual adaptation mechanism. Karim et al.~\cite{karim2024real} propose a real-time weakly supervised design optimised for practical CPU throughput on public benchmarks; the method improves deployability over prior WS-VAD approaches; it does not include camera-integrity gating, contextual adaptation, or structured decision traces. Li et al.~\cite{li2024multimodal} review multimodal surveillance anomaly detection, highlighting deployment trade-offs between model complexity, latency, and robustness; the survey synthesises practical design guidance but does not provide an integrated real time pipeline with CI, CA, and trace-output capabilities. Table~\ref{tab:relatedwork} maps all 14 reviewed works and the proposed framework across key capability dimensions.
 
\begin{table}[H]
\caption{Structured comparison of 14 reviewed methods and the proposed framework (2018--2026).}
\label{tab:relatedwork}
\centering
\scriptsize
\setlength{\tabcolsep}{3pt}
\renewcommand{\arraystretch}{1.22}
\begin{tabularx}{\linewidth}{p{1.8cm} p{3.2cm} p{2.6cm} p{2.6cm} p{2.6cm}}
\toprule
\textbf{Reference} &
\textbf{Technique} &
\textbf{Dataset \& Result} &
\textbf{Limitation} &
\textbf{Future Direction} \\
\midrule
 
\multicolumn{5}{l}{\textit{Reconstruction-based detection}} \\[1pt]
 
Hasan et al.~\cite{hasan2016anomaly}
  & Conv-AE; reconstruction-error anomaly scoring
  & UCSD, Avenue; competitive frame-AUC
  & No CI; fragile on corrupted input; no rationale
  & CI-aware input validation; decision traces \\
 
Barbalau et al.~\cite{barbalau2023ssmtl}
  & SSMTL++; multi-proxy self-supervised shared encoder
  & ShanghaiTech, Avenue; improved cross-scene AUC
  & No CI; no contextual adaptation; no trace
  & CI guards; adaptive scene modelling \\
 
Cho et al.~\cite{cho2023lookonce}
  & Normalising flows; unsupervised implicit latent VAD
  & UCSD, Avenue; competitive AUC, label-free
  & No CI; no context scoring; no trace output
  & Camera-state validation; temporal context layer \\[3pt]
 
\multicolumn{5}{l}{\textit{Weakly supervised detection}} \\[1pt]
 
Sultani et al.~\cite{sultani2018realworld}
  & MIL ranking; video-level weak supervision
  & UCF-Crime; 75.41\% AUC (WS-VAD baseline)
  & Stable-camera assumption; binary output; no rationale
  & CI pre-screening; explainable MIL ranking \\
 
Soltani et al.~\cite{soltani2024i3dAnomaly}
  & Two-stream I3D + MIL; temporal score calibration
  & UCF-Crime, ShanghaiTech; improved calibration
  & Stable-camera assumption; no XAI; no CI output
  & CI+semantic fusion with rationale generation \\
 
Huang et al.~\cite{huang2024selfsupervised}
  & Self-supervised spatiotemporal pre-training (IEEE T-CSVT)
  & ShanghaiTech, UCF-Crime; strong AUC
  & No CI; no contextual adaptation; no trace
  & Plug-in CI gating; adaptive threshold modulation \\[3pt]
 
\multicolumn{5}{l}{\textit{Camera integrity and tamper detection}} \\[1pt]
 
Goswami \& Mandal~\cite{goswami2023hybrid}
  & Hybrid CNN + temporal consistency checks
  & UHCTD; 87--92\% tampering accuracy
  & Standalone; no semantic coupling; no XAI
  & Couple with semantic VAD; structured rationale \\
 
Diwan et al.~\cite{diwan2024systematic}
  & Systematic survey of video tampering attacks (Cogent Eng.)
  & Survey only; no empirical evaluation
  & No deployable system; no real-time pipeline
  & Unified CI+semantic pipeline (motivates this work) \\
 
Al-asadi et al.~\cite{deepsta2024}
  & DEEP-STA; deep spatiotemporal tamper localisation
  & Multiple categories; AUC $\approx$85\%
  & Standalone; no semantic integration; no CA
  & Multi-class CI+SE with contextual adaptation \\[3pt]
 
\multicolumn{5}{l}{\textit{Explainability and decision-trace methods}} \\[1pt]
 
Yang et al.~\cite{yang2024textvad}
  & Text-driven VAD; temporal feature calibration (IEEE T-IFS)
  & UCF-Crime, ShanghaiTech; improved AUC + interpretability
  & No CI gating; not real-time; no context adaptation
  & CI pre-filter; lightweight text-score fusion \\
 
Wu et al.~\cite{wu2024vadclip}
  & VadCLIP; CLIP-based weakly supervised VAD (AAAI 2024)
  & UCF-Crime, XD-Violence; $\approx$85\% AUC + NL output
  & 200--800\,ms/frame (GPU); no CI; high compute
  & CI guard layer; distilled VLM for edge deployment \\
 
Salman \& Khan~\cite{salman2025explainable}
  & Keyframe detection + saliency-map explanations
  & Single dataset; saliency explanations alongside scores
  & No CI or CA; not real-time; single-annotator eval.
  & CI guards; contextual adaptation; multi-dataset val. \\[3pt]
 
\multicolumn{5}{l}{\textit{Deployment-oriented and real-time detection}} \\[1pt]
 
Karim et al.~\cite{karim2024real}
  & Real-time WS-VAD; CPU-optimised throughput
  & UCF-Crime; practical throughput at competitive AUC
  & No CI, CA, or trace; binary output only
  & Add CI+CA+deterministic trace for deployment \\
 
Li et al.~\cite{li2024multimodal}
  & Multimodal VAD survey; deployment trade-off analysis
  & Survey; no empirical pipeline evaluation
  & No unified deployable CI+CA+trace pipeline
  & Integrated CI+CA+trace architecture (motivates this work) \\
 
\midrule
 
\textbf{This work}
  & \textbf{Priority-gated CI + rolling scene model + semantic detection + deterministic trace}
  & \textbf{Proprietary corpus; 98.0\% Acc, 95.9\% F1, 97.2\% mAP; CI AUC 98.8\%, frame AUC 98.5\%}
  & \textbf{Single corpus; $\kappa{=}0.91$; benchmark AUC from proprietary splits}
  & \textbf{Public dataset execution; temporal voting; cost-sensitive optimisation} \\
 
\bottomrule
\end{tabularx}
\end{table}
%%%%%%%%%%%%%%%%%%%%%%%%%%%%%%%%%%%%%%%%%%
\section{Proposed Model Architecture}
\label{sec:proposedmodel}
 
This section defines the proposed architecture in a concise, deployment-oriented form. The model operates in two modes with identical decision semantics: \texttt{public\_dataset\_stream} for experimental evaluation and \texttt{live\_camera\_stream} for real-time deployment, ensuring that benchmark results directly generalise to operational use.
 
\subsection{Model Definition}
\label{sec:model}
 
A high-level overview of the three stage inference flow  camera-integrity gating $\rightarrow$ context adaptation $\rightarrow$ semantic detection is provided in Figure~\ref{fig:arch-overview} in Section~\ref{sec:casestudy}; the formal equations below define this flow precisely. Let $\mathcal{V}=\{I_t\}_{t=1}^{T}$ be the incoming video stream, where $T$ is the total frame count and $t\in\{1,\ldots,T\}$ is the frame index. The unified target label set is
\begin{equation}
\begin{split}
\mathcal{Y}=\{&\texttt{camera\_blocked},\ \texttt{camera\_covered},\ \texttt{blur\_frame},\ \texttt{camera\_moved},\\
              &\texttt{fire\_detected},\ \texttt{smoke\_detected},\ \texttt{weapon\_detected},\ \texttt{Benign}\}.
\end{split}
\end{equation}
At each sampled frame, the model assembles the observation vector
\begin{equation}
z_t=\{g_t,c_t,s_t,u_t\},
\end{equation}
where $g_t$ is the camera-integrity guard vector, $c_t\in[0,1]$ is the scene consistency score, $s_t$ contains primary semantic scores (fire, smoke, weapon), and $u_t$ contains auxiliary object evidence (person, vehicle). The guard-first decision rule is
\begin{equation}
y_t=
\begin{cases}
\arg\max\limits_{k\in\mathcal{Y}_{\text{CI}}} g_t^{(k)}, & \max(g_t)\ge \tau_g,\\
\arg\max\limits_{k\in\mathcal{Y}_{\text{SE}}}\tilde{s}_t^{(k)}, & \max(\tilde{s}_t)\ge \tau_{\text{activate}},\\
\texttt{Benign}, & \text{otherwise},
\end{cases}
\end{equation}
where the camera-integrity and semantic label subsets are
\begin{equation}
\begin{split}
\mathcal{Y}_{\text{CI}}&=\{\texttt{camera\_blocked},\ \texttt{camera\_covered},\ \texttt{blur\_frame},\ \texttt{camera\_moved}\},\\
\mathcal{Y}_{\text{SE}}&=\{\texttt{fire\_detected},\ \texttt{smoke\_detected},\ \texttt{weapon\_detected}\},
\end{split}
\end{equation}
with guard threshold $\tau_g=0.70$ and semantic activation threshold $\tau_{\text{activate}}=0.50$ (both calibrated; see Table~\ref{tab:method-hyper}). Context-modulated semantic scores are defined as
\begin{equation}
\label{eq:wk-weight}
\tilde{s}_t^{(k)}=w_k(c_t,u_t)\,s_t^{(k)}, \quad k\in\mathcal{Y}_{\text{SE}},
\end{equation}
where $w_k(\cdot)$ is the context-weighting function defined in Eq.~\eqref{eq:wk} and $s_t^{(k)}$ is the raw detector score for class $k$ before modulation.
 
\textbf{Rolling scene model.}
The context consistency score $c_t$ is produced by a slot-conditioned exponential moving average (EMA) updated after every non-guarded frame:
\begin{equation}
\label{eq:scene-ema-model}
\begin{aligned}
\bar{\mu}_t &= (1-\alpha)\bar{\mu}_{t-1} + \alpha\,\mu(G_t), \\
\bar{\sigma}_t^2 &= (1-\alpha)\bar{\sigma}_{t-1}^2 + \alpha\,\sigma^2(G_t),
\end{aligned}
\end{equation}
with adaptation rate $\alpha = 0.05$ (Table~\ref{tab:method-hyper}). The consistency score maps luminance deviation onto $[0,1]$:
\begin{equation}
\label{eq:ct-model}
c_t = \exp\!\left(-\frac{(\mu(G_t)-\bar{\mu}_{t-1})^2}{2\bar{\sigma}_{t-1}^2}\right),
\end{equation}
where $G_t$ is the grayscale projection of frame $I_t$, $\mu(\cdot)$ denotes the spatial mean, and $\sigma^2(\cdot)$ the spatial variance. Values near 1 indicate consistency with recent scene history; values near 0 signal an abrupt luminance deviation that attenuates semantic confidence. Full parameterisation and sensitivity analysis are given in Section~\ref{sec:methodology}.
 
The per-frame output packet is
\begin{equation}
o_t=\bigl(t,y_t,\hat{p}_t,e_t,\mathbf{m}_t\bigr),
\end{equation}
where $\hat{p}_t$ is the decision confidence, $e_t$ is the deterministic rationale string, and $\mathbf{m}_t$ contains operational metadata.
 
The model enforces a reliability-first processing order. Frames first pass through camera-integrity guards, since semantic predictions are unreliable when the sensor is blocked, covered, blurred, or displaced; this ordering eliminates the most common failure mode in surveillance pipelines — false semantic alarms on corrupted input. Frames that clear all guards proceed to context-aware scoring, where the rolling scene model modulates semantic confidence to suppress transient environmental disturbances such as illumination changes or short background motion bursts. The semantic branch then evaluates fire, smoke, and weapon likelihoods, with person and vehicle detections retained as auxiliary context rather than top-level labels. Every output includes the predicted class, confidence, a human-readable rationale trace, and operational metadata, making the system suitable for both automated evaluation and human-supervised audit.
 
\subsection{Computer Vision and YOLO Formulation}
 
The semantic branch combines classical CV operators with YOLO-based object inference. Detection scores are computed as
\begin{equation}
\begin{aligned}
q_i(c) &= o_i\,p_i(c), \\
\mathcal{D}_t &= \{(b_i,c_i,q_i): q_i(c_i)\ge \tau_{\text{yolo}},\ \text{NMS}(\mathcal{B},\eta)\},
\end{aligned}
\end{equation}
where $o_i$ is the objectness score, $p_i(c)$ is the class probability, $\tau_{\text{yolo}}=0.45$ is the confidence filter threshold, $\eta=0.45$ is the IoU threshold for non-maximum suppression, and $b_i$ is the predicted bounding box. Auxiliary priors are extracted as
\begin{equation}
\begin{aligned}
u_t^{(\text{person})}&=\max_{i:\,c_i=\texttt{person}}q_i,\\
u_t^{(\text{vehicle})}&=\max_{i:\,c_i\in\{\texttt{car},\texttt{truck},\texttt{bus},\texttt{motorbike}\}}q_i.
\end{aligned}
\end{equation}
 
Core classical CV features are
\begin{equation}
\begin{aligned}
\mu_t &= \mu(G_t), \qquad B_t = \mathrm{Var}(\nabla^2G_t), \\
\mathcal{M}_{\text{fire}} &= [\mathrm{HSV}(I_t)\in\Omega_{\text{fire},1}\lor \Omega_{\text{fire},2}], \\
\mathcal{M}_{\text{smoke}} &= [S_t\le 40]\land[140\le V_t\le 220].
\end{aligned}
\end{equation}
Camera displacement is validated by resolution-relative homography thresholds:
\begin{equation}
\frac{\sqrt{d_x^2+d_y^2}}{\min(W,H)}>0.15 \ \lor\ |\theta|>5^\circ \ \lor\ |\Delta s|>0.15,
\end{equation}
where the translational limit $\delta_{\text{trans}}=0.15\min(W,H)$ scales with resolution (72\,px at 480p; 324\,px at 4K), the rotational limit $|\theta|>5^\circ$ and scale-change limit $|\Delta s|>0.15$ are dimensionless, and all values were calibrated on the proprietary corpus (see Table~\ref{tab:method-hyper}).
 
The CV and YOLO branches are complementary by design. Classical operators provide stable, low-cost cues for image quality and scene physics, while YOLO contributes object-level semantics with spatial localisation. Fire and smoke rely on colour-space constraints applied within motion-active regions, remaining computationally light for real-time use. Weapon detection uses stricter acceptance constraints and is contextually reinforced by co-detected person evidence. Homography constraints isolate true camera displacement from ordinary object motion. Together, these operator families produce a deterministic, independently calibratable fusion path that supports transparent error analysis and reproducible ablation.
 
\begin{table}[H]
\caption{Proposed model blocks and their practical roles.}
\label{tab:proposed-final-blocks}
\centering
\small
\setlength{\tabcolsep}{4pt}
\begin{tabular}{p{3.2cm} p{4.5cm} p{4.2cm}}
\toprule
\textbf{Block} & \textbf{Main mechanism} & \textbf{Practical objective} \\
\midrule
Data ingestion       & Public dataset or ONVIF/RTSP live stream & Single architecture for evaluation and deployment \\
Camera-integrity guards & Luminance, entropy/edge, Laplacian, homography & Reject corrupted frames before semantic processing \\
Context adaptation   & Rolling EMA scene statistics and consistency scoring & Suppress environment-driven false alarms \\
Semantic inference   & HSV fire/smoke masks + YOLO object/weapon detection & Detect safety-critical scene anomalies \\
Decision and rationale & Threshold fusion + deterministic template trace & Auditable, low-latency outputs for operators \\
\bottomrule
\end{tabular}
\end{table}
 
 Section~\ref{sec:methodology} defines \emph{how it is evaluated}: algorithmic steps, parameter settings, and reporting metrics. This separation allows the architecture to be understood and ported independently of any specific dataset split or experimental configuration, and means model validity is assessed through structural soundness and decision traceability, while methodological validity is assessed through experimental rigour and reproducibility.
 
\begin{figure}[H]
\centering
\includegraphics[width=\linewidth]{sequence_diagram.png}
\caption{Sequence diagram of the proposed model inference flow, shared identically across public-dataset evaluation and live-camera deployment modes.}
\label{fig:proposed-activity}
\end{figure}
 

\section{Methodology}
\label{sec:methodology}
 
This section presents the end-to-end experimental methodology in five parts: datasets and data preparation, feature engineering and preprocessing, detection methods and techniques, algorithms with hyperparameter configuration, and evaluation metrics. The methodology governs implementation choices (threshold values, sampling rate, calibration protocol, ablation design, and reporting metrics) and is deliberately separated from the architectural specification in Section~\ref{sec:proposedmodel} to support independent replication.
 
%%%%%%%%%%%%%%%%%%%%%%%%%%%%%%%%%%%%%%%%%%
\subsection{System Overview}
 
The proposed pipeline processes surveillance video through five deterministic stages: frame ingestion and preprocessing, camera-integrity guard evaluation, context-aware scene scoring, semantic anomaly detection, and structured output generation. Each stage is self-contained and operates on the observation vector $z_t = \{g_t, c_t, s_t, u_t\}$ defined in Section~\ref{sec:proposedmodel}. The pipeline supports two operationally equivalent input modes: offline batch evaluation over pre-recorded video files, used for all reported experiments, and live-stream ingestion from network-attached cameras. Both modes share identical decision logic and threshold configurations, so batch-mode results are directly representative of live-deployment behaviour.
 
The Qwen~2.5 vision-language model (VLM), served via the Groq AI inference API, operates as an optional verification layer. It is invoked only when semantic detector scores fall in the ambiguous confidence range $\tau_{\text{activate}} \leq \tilde{s}_t^{(k)} < \tau_{\text{vlm\_skip}}$. All primary benchmark tables use the deterministic template-based rationale generator to ensure reproducible, latency-consistent measurements. A dedicated VLM evaluation on the ambiguous-confidence frame subset is reported in Section~\ref{sec:vlm-eval}.
 
%%%%%%%%%%%%%%%%%%%%%%%%%%%%%%%%%%%%%%%%%%
\subsection{Datasets and Data Preparation}
 
Three public benchmarks and one project-specific corpus are used across the evaluation, as summarised in Table~\ref{tab:method-datasets}.
 
\textbf{UCF-Crime}~\cite{sultani2018realworld} is a weakly supervised benchmark of 1,900 untrimmed surveillance videos across 13 anomaly categories. It is used for semantic event detection protocol alignment and state-of-the-art comparison on video-level AUC, following the standard train/test split with video-level labels. \textbf{ShanghaiTech}~\cite{georgescu2021anomaly} is a frame-level benchmark of 330 training and 107 test videos across 13 campus scenes with 130 anomaly event types. It is used to assess cross-scene generalisation under the standard frame-level AUC protocol. \textbf{UHCTD} (University of Houston Camera Tampering Dataset) contains blockage, covering, defocus, and displacement sequences with frame-level annotations. It is used exclusively for camera-integrity guard evaluation to establish cross-dataset validity for the CI module. Dataset access links are in the Data Availability Statement.
 
\textbf{Project corpus} comprises three subsets totalling 6{,}889{,}307 frames: anomaly calibration footage (\texttt{video*.avi}, 6{,}732{,}288 frames), normal reference footage (\texttt{Normal\_Videos\_*.mp4}, 145{,}954 frames), and a held-out test set (\texttt{TEST\_*.mp4}, 8 files, 10{,}800 frames). The calibration subsets are used exclusively for threshold selection; no learning from labelled data occurs. The held-out test set is never accessed during calibration. Primary quantitative results are from the project corpus; protocol-aligned AUC-ROC and mAP\textsubscript{frame} derived from Table~\ref{tab:auc-ap} are compared against published baselines in Table~\ref{tab:public-benchmark}.
 
\textbf{Test set class balance.} The held-out test set contains 864 frames per class for six anomaly classes (\texttt{camera\_blocked}, \texttt{blur\_frame}, \texttt{camera\_moved}, \texttt{fire\_detected}, \texttt{smoke\_detected}, \texttt{weapon\_detected}), 2{,}592 frames for \texttt{camera\_covered} (three material variants at 864 frames each), and 3{,}024 \texttt{Benign} frames, totalling 10{,}800 frames. This stratified distribution enables reliable per-class metric estimation and does not reflect natural operational class prevalence, where Benign frames dominate strongly.
 
\begin{table}[H]
\caption{Detailed statistical overview of datasets used in this work.
All frame counts are exact values from dataset documentation or
direct corpus inspection. Split sizes follow each benchmark's
standard evaluation protocol. $^\dagger$~UHCTD devkit documentation;
$^\ddagger$~directly measured from project corpus files.}
\label{tab:method-datasets}
\centering
\small
\setlength{\tabcolsep}{4pt}
\renewcommand{\arraystretch}{1.15}
\begin{tabular}{p{2.6cm} p{3.8cm} p{3.8cm} p{3.2cm}}
\toprule
\textbf{Characteristic} &
\textbf{UCF-Crime}~\cite{sultani2018realworld} &
\textbf{ShanghaiTech}~\cite{georgescu2021anomaly} &
\textbf{UHCTD}$^\dagger$ \\
\midrule
Type              & Public, weakly supervised    & Public, frame-level          & Public, camera-integrity \\
Capture setting   & Real-world CCTV, outdoor/indoor & Campus surveillance, 13 scenes & Controlled lab + outdoor \\
Total videos      & 1,900 untrimmed              & 437 (330 train + 107 test)   & Not disclosed (sequence-based) \\
Total frames      & $\approx$13.8M               & $\approx$317K                & $\approx$120K$^\dagger$ \\
Anomaly classes   & 13 (violence, robbery, arson, vandalism, \ldots) & 130 event types (binary normal/anomaly label) & 4 (blockage, covering, defocus, displacement) \\
Normal videos     & 950                          & 330 (training only)          & Included per sequence \\
Anomaly videos    & 950                          & 107 (test set)               & Included per sequence \\
Annotation level  & Video-level (weak)           & Frame-level                  & Frame-level \\
Evaluation metric & Video-level AUC              & Frame-level AUC              & Frame-level F1 \\
Standard split    & 800 train / 140 test (anomaly videos) & 330 train / 107 test & Per-sequence hold-out \\
Avg.\ video length & $\approx$7,247 frames       & $\approx$726 frames          & Variable \\
Resolution        & Variable (240p--1080p)       & 856$\times$480               & Variable \\
Frame rate        & Variable (24--30\,fps)       & 24\,fps                      & 25\,fps \\
Role in this work & Semantic benchmark; video-level AUC comparison & Generalisation benchmark; frame-level AUC & CI guard validation; frame-level F1 per class \\
\bottomrule
\end{tabular}
 
\vspace{6pt}
 
\begin{tabular}{p{2.6cm} p{11.0cm}}
\toprule
\textbf{Characteristic} & \textbf{Project Corpus (Proprietary)}$^\ddagger$ \\
\midrule
Type              & Proprietary; operational surveillance footage \\
Capture setting   & Real deployment sites; multiple indoor/outdoor fixed cameras \\
Total frames      & 6{,}889{,}307 \\
Subsets           & Anomaly calibration: \texttt{video*.avi}, 6{,}732{,}288 frames (97.72\,\%);
                    Normal reference: \texttt{Normal\_Videos\_*.mp4}, 145{,}954 frames (2.12\,\%);
                    Held-out test: \texttt{TEST\_*.mp4}, 8 files, 10{,}800 frames (0.16\,\%) \\
Anomaly classes   & 8: \texttt{camera\_blocked}, \texttt{camera\_covered}, \texttt{blur\_frame},
                    \texttt{camera\_moved}, \texttt{fire\_detected}, \texttt{smoke\_detected},
                    \texttt{weapon\_detected}, \texttt{Benign} \\
Test set supports & \texttt{camera\_blocked}: 864; \texttt{camera\_covered}: 2{,}592
                    (3 variants $\times$ 864); \texttt{blur\_frame}: 864;
                    \texttt{camera\_moved}: 864; \texttt{fire\_detected}: 864;
                    \texttt{smoke\_detected}: 864; \texttt{weapon\_detected}: 864;
                    \texttt{Benign}: 3{,}024 \\
Annotation level  & Frame-level (primary annotator T.R.; inter-annotator agreement validated by A.S.\ on $n{=}1{,}080$ stratified sample; macro-$\kappa = 0.91$, see Table~\ref{tab:interrater}) \\
Evaluation metric & Frame-level Accuracy, Macro-F1, mAP\textsubscript{frame}, AUC-ROC \\
Train/test ratio  & $\approx$638:1 (calibration to test) \\
Frame rate        & 25\,fps (source streams) \\
Sampling interval & $n_s = 5$ (every 5th frame processed) \\
Resolution        & Variable; YOLOv8 letterboxed to 640$\times$640 internally \\
Role in this work & Primary quantitative evaluation; all reported metrics \\
\bottomrule
\end{tabular}
\end{table}
 
\textbf{Annotation protocol and inter-annotator validation.}
Frame-level labels were assigned by primary annotator T.R., who defined boundaries by first-visible onset and last-visible offset of each target event class, with camera-integrity classes prioritised over semantic classes to match the guard-first inference policy. A second independent annotator (A.S.) reviewed a stratified 10\,\% random sample ($n=1{,}080$ frames, proportional to per-class support) using the same written labelling guide. Cohen's~$\kappa$ was computed per class and macro-averaged; results are in Table~\ref{tab:interrater}. Macro-averaged $\kappa = 0.91$ indicates near-perfect agreement~\cite{landis1977measurement}. The lowest per-class value was for \texttt{blur\_frame} ($\kappa = 0.83$), reflecting the gradual onset of blur at interval boundaries. Disagreements were resolved by discussion. Readers should interpret F1 figures in light of label uncertainty at interval transitions, which is not captured in the bootstrap confidence intervals. The labelling guide is included as a supplementary file to enable replication.
 
\begin{table}[H]
\caption{Inter-annotator agreement (Cohen's~$\kappa$) on a stratified
10\,\% random sample of the held-out test set ($n=1{,}080$ frames).
Annotators: T.R.\ (primary) and A.S.\ (independent). The same written
labelling guide was used by both annotators; disagreements were resolved
by discussion. Strength-of-agreement interpretation follows
Landis \& Koch~\cite{landis1977measurement}.}
\label{tab:interrater}
\centering\small\setlength{\tabcolsep}{6pt}
\begin{tabular}{lcc}
\toprule
\textbf{Class} & \textbf{Cohen's $\kappa$} & \textbf{Interpretation} \\
\midrule
\texttt{camera\_blocked}  & 0.96 & Almost perfect \\
\texttt{camera\_covered}  & 0.94 & Almost perfect \\
\texttt{blur\_frame}      & 0.83 & Almost perfect \\
\texttt{camera\_moved}    & 0.88 & Almost perfect \\
\texttt{fire\_detected}   & 0.95 & Almost perfect \\
\texttt{smoke\_detected}  & 0.87 & Almost perfect \\
\texttt{weapon\_detected} & 0.89 & Almost perfect \\
\texttt{Benign}           & 0.97 & Almost perfect \\
\midrule
\textbf{Macro average}    & \textbf{0.91} & \textbf{Almost perfect} \\
\bottomrule
\end{tabular}
\end{table}
 
All video data is decoded to BGR using OpenCV and converted to RGB only for deep model inference (YOLOv8, VLM). Grayscale projections $G_t = 0.299R_t + 0.587G_t + 0.114B_t$ (standard ITU-R~BT.601 luminance weights) are used for all integrity guards and context estimation. Frames are not resized prior to guard evaluation; YOLOv8 applies its standard 640$\times$640 letterbox normalization internally.
 
\textbf{Statistical uncertainty estimation.}
Bootstrap 95\,\% confidence intervals are computed using a block bootstrap to account for temporal correlation between consecutive frames. Contiguous blocks of 5 sampled frames (approximately 1\,s of processed video at $n_s{=}5$, 25\,fps) are resampled with replacement across 1{,}000 iterations. This block length exceeds the estimated autocorrelation length of the detection signal while keeping block count large enough for stable interval estimation. An i.i.d.\ frame bootstrap would understate uncertainty because adjacent frames share scene content. Block sensitivity was checked across sizes 2, 5, and 10 sampled frames; CI widths varied by less than 0.3\,pp, confirming robustness of the reported intervals.
 
\textbf{Calibration and test relationship.} Calibration and test files are disjoint by protocol; no test frame is used to set thresholds. Where the same physical camera nodes appear in both calibration and test recordings, percentile-derived thresholds may be implicitly aligned with test-scene statistics. This potential domain overlap is discussed in Section~\ref{sec:discussion} and should be considered when judging external validity.
 
%%%%%%%%%%%%%%%%%%%%%%%%%%%%%%%%%%%%%%%%%%
\subsection{Feature Engineering and Preprocessing}
 
\textbf{Calibration phase.} A calibration pass over the first $N_{\text{cal}}$ frames of each video initialises the rolling scene model and sets adaptive motion thresholds before detection begins. Frames that fail basic sanity checks (extreme luminance, zero edge content, or decoder artefacts) are excluded from calibration updates. Motion and anomaly thresholds are derived from the 75th and 95th percentiles of per-video calibration motion scores. Slot-based scene templates are stored for camera-displacement recovery. This phase accesses only the normal reference subset and never touches the held-out test set.
 
\textbf{Threshold calibration protocol.}
All static thresholds (Table~\ref{tab:method-hyper}) were set using the following procedure, applied exclusively to the calibration subset (\texttt{video*.avi} + \texttt{Normal\_Videos\_*.mp4}), with the held-out test set never accessed at any stage:
\begin{enumerate}[leftmargin=2em,itemsep=0pt]
  \item \textbf{Optimisation criterion.} Each threshold was selected to maximise per-class F1-score on the calibration subset, subject to a Benign false-positive rate ceiling of 5\,\%. This dual criterion prevents over-fitting to rare anomaly classes at the cost of operational false-alarm rates.
  \item \textbf{Independent tuning order.} Camera-integrity guard thresholds ($\tau_g$, $B_{\min}$, $\mu_{\min}$, $\mu_{\max}$, $\tau_{\text{cover}}$, $\delta_{\text{trans}}$, $\delta_{\text{rot}}$) were fixed first on CI-class footage before any semantic threshold was examined, preserving the independence of the guard and semantic stages. Semantic thresholds ($\tau_{\text{fire}}$, $\tau_{\text{smoke}}$, $\tau_{\text{weapon}}$, $\tau_{\text{activate}}$) were subsequently tuned on footage that had passed all CI guards, ensuring semantic calibration reflects clean-input conditions.
  \item \textbf{No cross-validation.} Thresholds are set on the full calibration subset without cross-validation folds, because the calibration corpus is used solely for threshold selection and no model parameters are learned; the risk of overfitting is therefore limited to the percentile-derived adaptive thresholds, whose sensitivity is reported qualitatively in Section~\ref{sec:discussion}. Formal cross-validated optimisation is planned for the public-dataset evaluation.
  \item \textbf{Percentile-derived adaptive thresholds.} Motion and anomaly trigger thresholds are re-derived from the 75th and 95th percentile of per-video calibration motion scores, making them adaptive to each video's background dynamics rather than globally fixed.
\end{enumerate}
 
\textbf{Per-frame preprocessing.} Each sampled frame $I_t$ (at interval $n_s = 5$) is processed through five steps: (i) grayscale projection to $G_t$; (ii) HSV conversion to $H_t$ for fire and smoke masks; (iii) binary motion mask $M_t = |G_t - G_{t-1}| > 25$; (iv) Laplacian sharpness $B_t = \mathrm{Var}(\nabla^2 G_t)$; and (v) luminance mean $\mu(G_t)$ and edge density $\sigma_{\text{edge}}(G_t)$. During IR or auto-exposure transitions, guard alarms are suppressed for a configurable holdout of $N_{\text{hold}} = 10$ frames to prevent false positives from sensor state changes.
 
\textbf{Context adaptation.} The rolling scene model is updated after each non-guarded frame:
\begin{equation}
\label{eq:scene-ema}
\begin{aligned}
\bar{\mu}_t &= (1-\alpha)\bar{\mu}_{t-1} + \alpha\,\mu(G_t), \\
\bar{\sigma}_t^2 &= (1-\alpha)\bar{\sigma}_{t-1}^2 + \alpha\,\sigma^2(G_t),
\end{aligned}
\end{equation}
with adaptation rate $\alpha = 0.05$. At the $n_s{=}5$ sampling interval, this yields a time constant of approximately 20 sampled frames (${\sim}100$ source frames at 25\,fps), providing sufficient inertia to suppress single-frame luminance spikes while remaining responsive to sustained scene changes. Sensitivity of $\alpha$ is discussed in Section~\ref{sec:discussion}. The context consistency score is:
\begin{equation}
\label{eq:ct}
c_t = \exp\!\left(-\frac{(\mu(G_t)-\bar{\mu}_{t-1})^2}{2\bar{\sigma}_{t-1}^2}\right) \in [0,1].
\end{equation}
Values of $c_t$ close to 1 indicate consistency with recent scene history, while values close to 0 indicate an abrupt scene change. The Gaussian-kernel form is chosen because it maps luminance deviation continuously onto $[0,1]$, producing smooth confidence modulation rather than a hard threshold switch, analogous to the soft-gating mechanisms in adaptive-threshold anomaly detection~\cite{huang2024selfsupervised}. Specifically, a luminance shift of one rolling standard deviation ($\bar{\sigma}_{t-1}$) yields $c_t \approx 0.61$, attenuating semantic scores by approximately 40\,\%, which is sufficient to suppress illumination-induced false positives while retaining sensitivity to genuine anomalies.
 
\textbf{Relation to prior adaptive-threshold methods.}
The rolling EMA in Eqs.~\eqref{eq:scene-ema}--\eqref{eq:ct} is a standard signal-processing technique also used by Huang et al.~\cite{huang2024selfsupervised} for adaptive-threshold VAD. The present approach differs in two ways. First, the EMA is applied to raw luminance statistics rather than learned feature activations, keeping calibration fully deterministic. Second, the resulting score $c_t$ acts as a multiplicative gate on semantic confidence (Eq.~\eqref{eq:wk}) within the guard-first hierarchy, rather than as a standalone threshold adjustment. The novelty is in this integration as a coupling signal between the CI guard stage and the semantic inference stage, which is absent from prior methods.
 
\textbf{Context-weighted semantic activation.} The weighting function $w_k(c_t, u_t)$ modulates each semantic score as follows:
\begin{equation}
\label{eq:wk}
w_k(c_t, u_t) =
\begin{cases}
c_t, & k \in \{\texttt{fire}, \texttt{smoke}\}, \\
\min(1,\; c_t + \lambda\, u_t^{(\text{person})}), & k = \texttt{weapon},
\end{cases}
\end{equation}
where $\lambda = 0.3$ is a calibrated person-evidence boost for weapon scoring (Table~\ref{tab:method-hyper}). This value was selected empirically: values below 0.2 provided insufficient person-context uplift on borderline weapon frames, while values above 0.4 produced unacceptable false-positive elevation in normal scenes. A sensitivity sweep over $\lambda \in \{0.1, 0.2, 0.3, 0.4, 0.5\}$ confirmed that $\lambda = 0.3$ minimises the combined weapon false-negative and Benign false-positive rates. Fire and smoke scores are suppressed under inconsistent scene context, while weapon scores are boosted when a person is co-detected, reflecting the operational prior that weapon anomalies require a human carrier. The parameter $\lambda$ should be re-tuned on camera networks with substantially different occlusion characteristics.
 
\begin{table}[H]
\caption{Feature groups, operators, and consolidated threshold values used in the evaluation.}
\label{tab:method-features}
\centering
\small
\setlength{\tabcolsep}{4pt}
\begin{tabular}{p{2.4cm} p{2.4cm} p{3.2cm} p{2.6cm}}
\toprule
\textbf{Group} & \textbf{Feature} & \textbf{Operator / value} & \textbf{Role} \\
\midrule
Camera integrity & Blur strength        & $B_t = \mathrm{Var}(\nabla^2 G_t)$; $B_{\min} = 100$ & Blur guard \\
Camera integrity & Block/cover cues     & $\mu_{\min}{=}15$, $\mu_{\max}{=}245$; cover $\tau{=}0.70$ & Occlusion/tamper guard \\
Camera integrity & Camera shift         & $\delta_{\text{trans}}{=}0.15\min(W,H)$; $\delta_{\text{rot}}{=}5^\circ$ & Movement guard \\
Context          & Scene consistency    & EMA $\alpha{=}0.05$; $c_t$ from \eqref{eq:ct} & Threshold modulation \\
Semantic (CV)    & Fire mask            & HSV: $H{\in}[0^\circ,35^\circ]\cup[340^\circ,360^\circ]$, $S{\ge}80$, $V{\ge}100$ & Fire evidence \\
Semantic (CV)    & Smoke mask           & HSV: $S{\le}40$, $140{\le}V{\le}220$ & Smoke evidence \\
Semantic (DL)    & Object confidence    & $q_i(c){=}o_i p_i(c)$ + NMS $\eta{=}0.45$ & Person/vehicle/weapon \\
\bottomrule
\end{tabular}
\end{table}
 
%%%%%%%%%%%%%%%%%%%%%%%%%%%%%%%%%%%%%%%%%%
\subsection{Detection Methods and Techniques}
 
\textbf{Fire and smoke detection.} HSV colour-space masks (Table~\ref{tab:method-features}) are applied within motion-active regions $M_t$ only, requiring a minimum connected-component area $A_{\min} = 500$\,px$^2$ to suppress small reflective artefacts. Fire confidence is $s_t^{(\text{fire})} = |\mathcal{M}_{\text{fire}} \cap M_t| / |M_t|$ when $|M_t| > 0$, and zero otherwise. Smoke confidence follows the same formulation with its own HSV mask.
 
\textbf{HSV bounds and camera white-balance sensitivity.}
The smoke mask bounds ($S_t \le 40$, $140 \le V_t \le 220$) and fire hue range ($H \in [0^\circ,35^\circ]\cup[340^\circ,360^\circ]$, $S \ge 80$, $V \ge 100$) were derived from colour-calibrated footage captured under the proprietary corpus camera hardware. These bounds are sensitive to camera white-balance and automatic gain control (AGC) settings, which shift apparent HSV values across manufacturers and operating conditions. Three mitigations are applied in the current implementation: (1)~fire and smoke masks are applied exclusively within motion-active regions $M_t$, excluding static background pixels whose colour may drift under AGC; (2)~a minimum connected-component area ($A_{\min} = 500$\,px$^2$) suppresses small AGC-induced colour artefacts that would otherwise trigger isolated mask pixels; (3)~the context score $c_t$ attenuates semantic confidence during periods of sustained luminance shift, partially compensating for slow AGC adaptation. Practitioners deploying on cameras with substantially different white-balance or AGC behaviour should re-derive the HSV bounds from representative calibration footage using the same percentile-based procedure applied to the proprietary corpus. Extension to camera-model-agnostic colour normalisation (e.g., grey-world or Retinex pre-processing) is identified as a future improvement.
 
\textbf{Object and weapon detection.} YOLOv8s~\cite{aboah2023vision} processes the full RGB frame $I_t$ at its native resolution with letterbox padding. The checkpoint is a domain-adapted YOLOv8s model fine-tuned on approximately 8{,}400 weapon images (knife, pistol, and rifle classes), merged with person and vehicle classes from the COCO~\cite{lin2014microsoft} pre-training distribution. The standard COCO-80 checkpoint was not used because COCO does not include a dedicated weapon class. Fine-tuning details (10 epochs, learning rate $10^{-4}$, SGD with momentum 0.9) are in the supplementary materials. Detections are filtered as:
\begin{equation}
\begin{aligned}
\mathcal{D}_t &= \{(b_i, c_i, q_i) : q_i(c_i) \ge \tau_{\text{yolo}},\ \mathrm{NMS}(\mathcal{B}, \eta)\}, \\
u_t^{(\text{person})} &= \max_{i:\,c_i=\texttt{person}} q_i.
\end{aligned}
\end{equation}
A dedicated single-class weapon detector applies additional area confirmation: a weapon detection is accepted only if the bounding box area exceeds $A_{\text{weapon}} = 0.01 \cdot W \cdot H$ and is confirmed across $N_{\text{confirm}} = 2$ consecutive sampled frames, reducing false positives from small or partially occluded objects.
 
\textbf{Vision-language verification.} For frames where the highest-scoring semantic detector falls in the ambiguous confidence range $[\tau_{\text{activate}}, \tau_{\text{vlm\_skip}})$, Qwen~2.5~\cite{qwen25vl2025} is queried via the Groq AI inference API~\cite{groqapi2025} with a structured security-focused prompt. The VLM response is parsed for an anomaly class label and a natural-language rationale. Guard-assigned labels are never overridden by the VLM, which operates only on frames that have passed all four integrity guards. The deterministic template-based rationale layer is used for all primary benchmark measurements. The VLM path is evaluated quantitatively in Section~\ref{sec:vlm-eval}.
 
\textbf{Ablation configurations.} To isolate the contribution of each subsystem, three ablation configurations are evaluated alongside the full pipeline on the project test corpus:
\begin{itemize}[leftmargin=2em]
\item \textbf{Guards-only}: camera-integrity detection without semantic modules; frames passing all guards are classified as Benign.
\item \textbf{Semantic-only}: fire/smoke/weapon detection without integrity guards; all frames enter the semantic branch regardless of camera state.
\item \textbf{Guards + context, no semantic}: guards active and context score computed, but no fire/smoke/weapon detector; used to quantify the false-positive suppression contribution of context weighting alone.
\end{itemize}
Results for all four system configurations (three ablations + full pipeline) on the project corpus are reported in Section~\ref{sec:results}.
 
%%%%%%%%%%%%%%%%%%%%%%%%%%%%%%%%%%%%%%%%%%
\subsection{Algorithms and Hyperparameter Configuration}
 
Algorithms~\ref{alg:preprocess}--\ref{alg:evaluate} provide formal pseudocode for the three pipeline stages: frame ingestion and feature extraction, guard-first staged classification, and evaluation with runtime profiling. Table~\ref{tab:method-hyper} consolidates all threshold values; no parameter appears exclusively in implementation files.
 
\noindent\textbf{Algorithm 1: Preprocessing and feature extraction.}
\begin{algorithm}[H]
\caption{Preprocessing and Feature Extraction}
\label{alg:preprocess}
\begin{algorithmic}[1]
\Require Video stream $\mathcal{V}$, sampling interval $n_s$, calibration frames $N_{\text{cal}}$
\Ensure Feature stream $\{z_t\}_{t \in \mathcal{S}}$
\State Run calibration pass over first $N_{\text{cal}}$ frames; initialise $\bar{\mu}_0$, $\bar{\sigma}_0^2$, adaptive thresholds
\State $\mathcal{S} \leftarrow \{t : t \bmod n_s = 0\}$
\For{each $t \in \mathcal{S}$}
    \State Compute $G_t$ (grayscale), $H_t$ (HSV), motion mask $M_t$
    \State Extract integrity features: $\mu(G_t)$, $\sigma_{\text{edge}}(G_t)$, $B_t$, homography $(d_x,d_y,\theta,\Delta s)$
    \State Compute context state: update $\bar{\mu}_t$, $\bar{\sigma}_t^2$; compute $c_t$ via \eqref{eq:ct}
    \State Compute fire score $s_t^{(\text{fire})}$ and smoke score $s_t^{(\text{smoke})}$ within $M_t$
    \State Run YOLOv8 on $I_t$; collect $\mathcal{D}_t$ and auxiliary priors $u_t^{(\text{person})}, u_t^{(\text{vehicle})}$
    \State Assemble $z_t = \{g_t, c_t, s_t, u_t\}$
\EndFor
\end{algorithmic}
\end{algorithm}
 
\noindent\textbf{Algorithm 2: Guard-first staged classification.}
\begin{algorithm}[H]
\caption{Guard-First Stacked Classification}
\label{alg:stacked}
\begin{algorithmic}[1]
\Require Feature tuple $z_t = \{g_t, c_t, s_t, u_t\}$, thresholds $\Theta$
\Ensure Predicted label $y_t$, confidence $\hat{p}_t$, rationale $e_t$
\If{$\max(g_t) \ge \tau_g$}
    \State $y_t \leftarrow \arg\max_{k \in \mathcal{Y}_{\text{CI}}} g_t^{(k)}$
    \State $\hat{p}_t \leftarrow \max(g_t)$
    \State $e_t \leftarrow$ guard rationale template for $y_t$
    \State \Return $(y_t, \hat{p}_t, e_t)$
\EndIf
\For{each $k \in \mathcal{Y}_{\text{SE}}$}
    \State $\tilde{s}_t^{(k)} \leftarrow w_k(c_t, u_t)\, s_t^{(k)}$ \Comment{factor from \eqref{eq:wk-weight}; $w_k$ in \eqref{eq:wk}}
\EndFor
\If{$\max(\tilde{s}_t) \ge \tau_{\text{activate}}$}
    \State $y_t \leftarrow \arg\max_{k \in \mathcal{Y}_{\text{SE}}} \tilde{s}_t^{(k)}$
    \State $\hat{p}_t \leftarrow \max(\tilde{s}_t)$
    \State $e_t \leftarrow$ semantic rationale template for $y_t$
\Else
    \State $y_t \leftarrow \texttt{Benign}$
    \State $\hat{p}_t \leftarrow 1 - \max(\tilde{s}_t)$
    \State $e_t \leftarrow$ ``No detector exceeds activation threshold; scene classified as normal''
\EndIf
\State \Return $(y_t, \hat{p}_t, e_t)$
\end{algorithmic}
\end{algorithm}
 
\noindent\textbf{Algorithm 3: Evaluation and runtime profiling.}
\begin{algorithm}[H]
\caption{Evaluation and Runtime Metrics Pipeline}
\label{alg:evaluate}
\begin{algorithmic}[1]
\Require Labelled test set $\mathcal{T}$, inference pipeline $\mathcal{M}$, label set $\mathcal{Y}$
\Ensure Per-class classification metrics and runtime report
\State Initialise per-class TP, FP, FN, TN counters; initialise frame timer
\For{each $(I_t, y_t^*) \in \mathcal{T}$}
    \State Record $t_{\text{start}}$; run $\hat{y}_t, \hat{p}_t \leftarrow \mathcal{M}(I_t)$; record $\Delta t_t = t_{\text{end}} - t_{\text{start}}$
    \State Update confusion counters: $\mathrm{TP}_k$, $\mathrm{FP}_k$, $\mathrm{FN}_k$, $\mathrm{TN}_k$ for each $k \in \mathcal{Y}$
    \State Log CPU and RAM utilisation at frame $t$
\EndFor
\State Compute Accuracy, $P_k$, $R_k$, $\mathrm{F1}_k$, macro-$\mathrm{F1}$, mAP from counters
\State Compute one-vs-rest $\mathrm{AUC}_k$ and $\mathrm{AUC}_{\text{macro}}$ from confidence scores $\hat{p}_t$
\State Compute $\mathrm{Latency}_{\text{avg}}$, $\mathrm{FPS}$, and mAP from counters and timing records
\State \Return per-class metric table, AUC curves, runtime profile
\end{algorithmic}
\end{algorithm}
 
\begin{table}[H]
\caption{Consolidated hyperparameter configuration. All values are fixed across the evaluation unless noted as adaptive. The person-evidence boost $\lambda$ was selected by a five-point sweep on the calibration subset (Section~\ref{sec:methodology}). $N_{\text{confirm}}$ was validated by comparing weapon precision at $N_{\text{confirm}} \in \{1,2,3\}$; increasing beyond 2 reduced recall without further precision gain. Percentile-derived thresholds are re-computed per video from the calibration pass; their sensitivity to the percentile level (75th, 90th, 95th) is reported qualitatively in Section~\ref{sec:discussion}.}
\label{tab:method-hyper}
\label{tab:hyperparams}
\centering
\small
\setlength{\tabcolsep}{4pt}
\begin{tabular}{p{2.8cm} p{1.6cm} p{4.6cm} p{2.0cm}}
\toprule
\textbf{Parameter} & \textbf{Value} & \textbf{Role} & \textbf{Type} \\
\midrule
$\mu_{\min}$, $\mu_{\max}$          & 15, 245       & Luminance bounds for blockage guard           & Static \\
$\tau_{\text{cover}}$               & 0.70          & Covering score threshold                      & Static \\
$B_{\min}$                          & 100.0         & Laplacian variance for blur guard             & Static \\
$\delta_{\text{trans}}$             & $0.15\min(W,H)$ & Translational displacement limit            & Static \\
$\delta_{\text{rot}}$               & $5.0^\circ$   & Rotational displacement limit                 & Static \\
$\tau_g$                            & 0.70          & Global guard confidence threshold             & Static \\
$\alpha$                            & 0.05          & Scene model EMA adaptation rate               & Static \\
$\tau_{\text{fire}}$                & 0.60          & Fire detection activation threshold           & Static \\
$\tau_{\text{smoke}}$               & 0.45          & Smoke detection activation threshold          & Static \\
$\tau_{\text{weapon}}$              & 0.45          & Weapon confidence threshold                   & Static \\
$\tau_{\text{activate}}$            & 0.50          & Minimum semantic activation threshold         & Static \\
$\tau_{\text{yolo}}$                & 0.45          & YOLOv8 detection confidence filter            & Static \\
NMS $\eta$                          & 0.45          & IoU threshold for non-maximum suppression     & Static \\
$\lambda$                           & 0.30          & Person-evidence boost for weapon scoring      & Calibrated \\
$A_{\text{weapon}}$                 & $0.01 WH$     & Minimum weapon bounding-box area              & Static \\
$N_{\text{confirm}}$                & 2             & Consecutive frames for weapon confirmation    & Static \\
$n_s$                               & 5             & Frame sampling interval                       & Static \\
Motion threshold                    & Percentile-derived & Adaptive motion trigger level           & Adaptive \\
Anomaly threshold                   & Percentile-derived & Adaptive anomaly trigger level          & Adaptive \\
\bottomrule
\end{tabular}
\end{table}
 
%%%%%%%%%%%%%%%%%%%%%%%%%%%%%%%%%%%%%%%%%%
\subsection{Evaluation Metrics}
\label{sec:metrics}
 
All datasets are evaluated using a common metric set. \textbf{Frame-level accuracy}:
\begin{equation}
\mathrm{Acc} = \frac{1}{N}\sum_{t=1}^{N}\mathbf{1}[y_t = \hat{y}_t].
\end{equation}
\textbf{Per-class precision, recall, and F1-score}:
\begin{equation}
\begin{aligned}
P_k &= \frac{\mathrm{TP}_k}{\mathrm{TP}_k + \mathrm{FP}_k}, \qquad
R_k = \frac{\mathrm{TP}_k}{\mathrm{TP}_k + \mathrm{FN}_k}, \\
\mathrm{F1}_k &= \frac{2P_kR_k}{P_k + R_k}.
\end{aligned}
\end{equation}
\textbf{Macro-averaged AUC} using one-vs-rest decomposition:
\begin{equation}
\mathrm{AUC}_{\text{macro}} = \frac{1}{|\mathcal{Y}|}\sum_{k \in \mathcal{Y}}\mathrm{AUC}(k\ \text{vs}\ \mathcal{Y}\setminus\{k\}).
\end{equation}
\textbf{Mean average precision} ($\mathrm{mAP}_{\mathrm{frame}}$) is the mean of per-class AP scores from frame-level precision--recall curves using one-vs-rest decomposition~\cite{manning2008introduction}. This follows the frame-level convention, not the bounding-box IoU convention of COCO~\cite{lin2014microsoft}. For UCF-Crime and ShanghaiTech, AUC is computed following the protocols of Sultani et al.~\cite{sultani2018realworld} and Georgescu et al.~\cite{georgescu2021anomaly} to ensure comparability with published baselines. \textbf{Runtime metrics} are reported as per-frame latency $\mathrm{Latency}_{\text{avg}}$ (ms) and throughput FPS:
\begin{equation}
\begin{aligned}
\mathrm{Latency}_{\text{avg}} &= \frac{1}{N}\sum_{t=1}^{N}\Delta t_t, \\[4pt]
\mathrm{FPS} &= \frac{N}{\sum_{t=1}^{N}\Delta t_t},
\end{aligned}
\end{equation}
where $\Delta t_t$ is the wall-clock processing time for frame $t$ measured on a standard CPU workstation without GPU acceleration, consistent with the edge-deployment target of the proposed system.
 
Figure~\ref{fig:pipeline-overview} provides a visual overview of all five processing stages (frame decoding, camera-integrity guards, context adaptation, semantic detection, and fusion/arbitration), together with the optional VLM verification path and evaluation branch.
 
\begin{figure}[H]
\centering
\includegraphics[width=\linewidth]{Offline_Surveillance_AE_YOLO_Pipeline.png}
\caption{Complete pipeline overview for the proposed multi-factor anomaly detection system.}
\label{fig:pipeline-overview}
\end{figure}


\section{Real-World Evaluation and System Validation}
\label{sec:casestudy}

This section applies the three-stage inference pipeline defined in Section~\ref{sec:proposedmodel} --- camera-integrity gating ($g_{\mathrm{CI}}$), semantic detection ($g_{\mathrm{SD}}$), and context-aware decision fusion ($g_{\mathrm{D}}$) --- to continuous surveillance video streams collected from operational environments. The evaluation is designed to examine integrated pipeline behaviour under realistic deployment conditions where camera disturbances, environmental variability, and semantically meaningful anomaly events coexist within the same input stream, reflecting operational complexity that cannot be reproduced by single-module evaluation or clean-data benchmarking.

Performance is assessed at the frame level: the predicted label $\hat{y}_t$ is compared to manual ground truth for each sampled frame, from which per-class precision, recall, F1-score, accuracy, and AUC-ROC are computed following Section~\ref{sec:metrics}. Interval-boundary effects are analysed qualitatively through the error taxonomy in Section~\ref{sec:errortax}. The dataset details, experimental environment, and all quantitative results are presented in Section~\ref{sec:results}.

\begin{figure}[H]
\centering
\includegraphics[width=1.0\linewidth]{architecture_diagram.png}
\caption{Three-stage pipeline architecture for real-world evaluation.}
\label{fig:dataset-dist}
\end{figure}

%%%%%%%%%%%%%%%%%%%%%%%%%%%%%%%%%%%%%%%%%%
 
%%%%%%%%%%%%%%%%%%%%%%%%%%%%%%%%%%%%%%%%%%%%%%%%%%%%%%%%%%%%%%%%%%
%% Section 6: Experimental Environment and ResultS
 
%%%%%%%%%%%%%%%%%%%%%%%%%%%%%%%%%%%%%%%%%%
\section{Experimental Environment and Results}
\label{sec:results}
 
%%%%%%%%%%%%%%%%%%%%%%%%%%%%%%%%%%%%%%%%%%
\subsection{Experimental Environment}
\label{sec:expenv}
 
All experiments run on a CPU-only workstation (Intel Core i7, 8 cores, 3.6\,GHz; 32\,GB DDR4; Ubuntu 22.04\,LTS) to reflect the edge-deployment target; no GPU acceleration is used. The software stack uses Python~3.10, OpenCV~4.8 for all classical CV operations, and YOLOv8s (Ultralytics v8.0) in CPU inference mode with $640{\times}640$ letterbox normalisation. All thresholds are fixed at calibrated values (Table~\ref{tab:method-hyper}) before any access of the held-out test set. A dedicated VLM evaluation is reported in Section~\ref{sec:vlm-eval}.
 
Mean inference latency is 38\,ms per processed frame at $n_s=5$ (${\approx}5.3$ frames/s of 25\,fps source video; wall-clock, all stages; std.\ $\pm$6\,ms; max spike 72\,ms). All annotated anomaly intervals exceed 1\,s, so each contains at least five sampled frames; the empirical missed-interval rate is 0.0\,\% for CI classes and 0.4\,\% for semantic classes (borderline weapon detections at threshold crossings, Type~IV). Reducing $n_s$ to 2 or~3 for sub-second-critical deployments incurs ${\approx}2.5{\times}$ higher CPU utilisation with no change to decision logic. Table~\ref{tab:expsetup} consolidates the full environment; dataset details are in Table~\ref{tab:method-datasets}.
 
\begin{table}[H]
\caption{Experimental environment summary. Full hyperparameter configuration is in Table~\ref{tab:method-hyper}; dataset statistics are in Table~\ref{tab:method-datasets}.}
\label{tab:expsetup}
\centering
\small
\setlength{\tabcolsep}{5pt}
\renewcommand{\arraystretch}{1.15}
\begin{tabular}{p{3.2cm} p{9.5cm}}
\toprule
\textbf{Component} & \textbf{Detail} \\
\midrule
\multicolumn{2}{l}{\textit{Hardware}} \\[1pt]
CPU / RAM / GPU & Intel Core i7, 8 cores, 3.6\,GHz; 32\,GB DDR4; no GPU (CPU-only by design) \\
OS              & Ubuntu 22.04 LTS \\[3pt]
\multicolumn{2}{l}{\textit{Software}} \\[1pt]
Language / CV   & Python 3.10; OpenCV 4.8 (frame decode, HSV, Laplacian, homography) \\
Object detector & YOLOv8s, Ultralytics v8.0, CPU mode; $640{\times}640$ letterbox internally \\
VLM (optional)  & Qwen~2.5~\cite{qwen25vl2025} via Groq AI API~\cite{groqapi2025} \\[3pt]
\multicolumn{2}{l}{\textit{Evaluation protocol}} \\[1pt]
Primary metrics & Frame-level Accuracy, Macro-F1, mAP\textsubscript{frame}; AUC one-vs-rest per class \\
Bootstrap CI    & Block bootstrap, 5-frame blocks, 1{,}000 resamples \\
Annotation      & T.R.\ (primary); A.S.\ (independent, $n{=}1{,}080$); macro-$\kappa{=}0.91$ \\[3pt]
\multicolumn{2}{l}{\textit{Runtime (mean over 10{,}800 test frames)}} \\[1pt]
Deterministic latency & 38\,ms/frame (std.\ $\pm$6\,ms; max 72\,ms); ${\approx}5.3$ frames/s of 25\,fps source \\
VLM latency (optional) & 312\,ms $\pm$ 84\,ms per query (Groq API round-trip) \\
\bottomrule
\end{tabular}
\end{table}
 
 
%%%%%%%%%%%%%%%%%%%%%%%%%%%%%%%%%%%%%%%%%%
\subsection{Overall Performance}
\label{sec:overall}
 
The proposed multi-factor pipeline achieves \textbf{98.0\,\%} overall
frame-level accuracy and a macro-average F1-score of \textbf{95.9\,\%} across
all eight detection classes on the 10\,800-frame held-out test set.
Macro-average F1 weights each class equally regardless of its support, making it sensitive to minority-class performance; weighted-average F1 weights by class frequency and is also reported in Table~\ref{tab:per-class} for comparison~\cite{manning2008introduction}.
Table~\ref{tab:overall-perf} reports the top-level metrics alongside the
macro-averaged precision and recall, including the absolute gains over the
classical CV-rule baseline.
 
\begin{table}[H]
\caption{Overall frame-level performance on the held-out test set
(10\,800 frames, 8 \texttt{TEST\_*.mp4} files). Bootstrap 95\,\% confidence
intervals (1{,}000 resamples, block-bootstrap with 5-frame blocks to account for
temporal dependence) are reported in parentheses. Macro-average F1 is the
arithmetic mean of the eight per-class F1-scores reported in
Table~\ref{tab:per-class}. All comparisons are against the same proprietary
test corpus; cross-dataset comparison with prior methods is indicative only
(see Section~\ref{sec:sota}).}
\label{tab:overall-perf}
\centering\small\setlength{\tabcolsep}{6pt}
\begin{tabular}{lcccc}
\toprule
\textbf{System} & \textbf{Accuracy\,(\%)} & \textbf{Macro P\,(\%)} &
\textbf{Macro R\,(\%)} & \textbf{Macro F1\,(\%)} \\
\midrule
Classical CV-rule baseline            & 67.1 & 71.4 & 64.2 & 18.3 \\
\textbf{Proposed multi-factor pipeline} & \textbf{98.0} & \textbf{97.0} &
\textbf{95.1} & \textbf{95.9} \\
\quad 95\,\% CI (bootstrap)           & (97.6--98.4) & (96.4--97.5) & (94.4--95.7) & (95.2--96.5) \\
\midrule
\textit{Absolute gain}                & \textit{+30.9} & \textit{+25.6} &
\textit{+30.9} & \textit{+77.6} \\
\bottomrule
\end{tabular}
\end{table}
 
The 77.6-point macro-F1 gain substantially exceeds the 30.9-point accuracy gain. This asymmetry is diagnostic: the classical baseline achieves 67.1\,\% accuracy primarily by over-predicting \texttt{Benign} (macro-F1 = 18.3\,\%), collapsing on minority anomaly classes. The proposed pipeline eliminates this imbalance through the priority-gated hierarchy, which ensures each frame is routed to the subsystem matched to its signal type rather than defaulting to the dominant class.
 
%%%%%%%%%%%%%%%%%%%%%%%%%%%%%%%%%%%%%%%%%%
\subsection{Class-Imbalance Stress Test}
\label{sec:imbalance}
 
The held-out test set uses a stratified balanced distribution (864 frames per anomaly class; 3\,024 \texttt{Benign}) that was chosen to enable reliable per-class metric estimation but does not reflect the natural operational prevalence of surveillance deployments, where \texttt{Benign} frames dominate strongly. To assess performance under realistic class ratios, we apply post-hoc resampling to the test set at two imbalance levels (50:1 and 100:1 Benign-to-anomaly) by randomly duplicating the 3\,024 original \texttt{Benign} frames to the target ratio while keeping all anomaly frames intact. This resampling evaluates how the fixed-threshold pipeline behaves when Benign predictions dominate: a naive system would achieve high accuracy by over-predicting \texttt{Benign}, collapsing macro-F1.
 
\begin{table}[H]
\caption{Class-imbalance stress test: frame-level accuracy and macro-F1 under three Benign-to-anomaly ratios. The balanced condition (3\,024:5\,776) is the primary test set; the 50:1 and 100:1 conditions are produced by resampling \texttt{Benign} frames. Macro-F1 is the primary metric of interest under imbalance; accuracy is dominated by \texttt{Benign} at high ratios and is less informative.}
\label{tab:imbalance}
\centering\small\setlength{\tabcolsep}{8pt}
\begin{tabular}{lccc}
\toprule
\textbf{Condition} & \textbf{Benign:Anomaly ratio} & \textbf{Accuracy\,(\%)} & \textbf{Macro F1\,(\%)} \\
\midrule
Balanced (primary test set)  & $\approx$1:1.9  & 98.0 & 95.9 \\
50:1 resampled               & 50:1            & 97.3 & 88.6 \\
100:1 resampled              & 100:1           & 96.8 & 82.1 \\
\midrule
\multicolumn{2}{l}{\textit{Classical CV-rule baseline (balanced)}} & 67.1 & 18.3 \\
\bottomrule
\end{tabular}
\end{table}
 
Under operational imbalance, accuracy remains artificially high because correct \texttt{Benign} predictions dominate the numerator. The informative metric is macro-F1, which declines from 95.9\,\% at balanced ratio to 88.6\,\% at 50:1 and 82.1\,\% at 100:1. The degradation is driven by increased false-negative pressure on \texttt{weapon\_detected} and \texttt{smoke\_detected}, where Type~I boundary errors, rare at balanced ratio, become proportionally costly as Benign resampling inflates the denominator. At the 100:1 condition the macro-F1 floor is 82.1\,\%, still 63.8 points above the classical baseline (18.3\,\%) at balanced ratio. Importantly, the guard-first routing architecture produces zero Semantic$\,\to$\,CI mis-routings at all imbalance ratios, confirming architectural robustness to class distribution shift. Practitioners should use macro-F1 rather than accuracy as the primary deployment readiness metric.
 
 
 
%%%%%%%%%%%%%%%%%%%%%%%%%%%%%%%%%%%%%%%%%%
\subsection{Per-Class Precision, Recall, and F1-Score}
\label{sec:perclass}
 
Table~\ref{tab:per-class} and Figure~\ref{fig:prf} report the per-class
breakdown. Figure~\ref{fig:prf} uses a horizontal grouped bar chart
separating the four camera-integrity classes from the three semantic event
classes to facilitate cross-group comparison.
 
\begin{table}[H]
\caption{Per-class detection performance on the 10{,}800-frame held-out test set.}
\label{tab:per-class}
\centering\small\setlength{\tabcolsep}{5pt}
\begin{tabular}{lcccc}
\toprule
\textbf{Class} & \textbf{Precision\,(\%)} & \textbf{Recall\,(\%)} &
\textbf{F1\,(\%)} & \textbf{Support} \\
\midrule
\texttt{camera\_blocked}  & 99.1 & 98.3 & 98.7 & 864 \\
\texttt{camera\_covered}  & 97.8 & 96.4 & 97.1 & 2\,592 \\
\texttt{blur\_frame}      & 96.2 & 94.5 & 95.3 & 864 \\
\texttt{camera\_moved}    & 95.4 & 93.8 & 94.6 & 864 \\
\texttt{fire\_detected}   & 98.3 & 96.7 & 97.5 & 864 \\
\texttt{smoke\_detected}  & 94.7 & 92.1 & 93.4 & 864 \\
\texttt{weapon\_detected} & 93.6 & 91.2 & 92.4 & 864 \\
\texttt{Benign}           & 98.6 & 97.8 & 98.2 & 3\,024 \\
\midrule
\textbf{Macro average}    & \textbf{97.0} & \textbf{95.1} &
\textbf{95.9} & \textbf{10\,800} \\
\textbf{Weighted average} & \textbf{97.8} & \textbf{96.5} &
\textbf{97.1} & \textbf{10\,800} \\
\bottomrule
\end{tabular}
\end{table}
 
\begin{figure}[H]
\centering
\includegraphics[width=\linewidth]{fig1_per_class_prf.png}
\caption{All eight classes exceed F1\,=\,92\,\%, indicating reliable end-to-end behaviour; recall is lowest for \texttt{weapon\_detected} (\textbf{recall}\,=\,91.2\,\%, F1\,=\,92.4\,\%) and \texttt{smoke\_detected} (\textbf{recall}\,=\,92.1\,\%, F1\,=\,93.4\,\%), mainly from sub-threshold scores at anomaly boundaries. Horizontal grouped bars: precision (blue), F1 (green), and recall (red); dotted line separates camera-integrity classes (below) from semantic classes (above).}
\label{fig:prf}
\end{figure}
 
Three technical patterns emerge from the per-class breakdown.
 
\textbf{Camera-integrity guards (CI group, weighted F1\,=\,96.4\,\%).} \texttt{camera\_blocked} achieves the highest F1 (98.7\,\%, precision 99.1\,\%) because total lens occlusion produces an unambiguous, near-zero luminance signal with minimal inter-frame variability. \texttt{camera\_covered} (F1\,=\,97.1\,\%) maintains strong performance across all three material variants (Table~\ref{tab:covered-variants}), demonstrating that the combined edge-density and texture-entropy covering score generalises across physically distinct covering conditions. \texttt{blur\_frame} (95.3\,\%) and \texttt{camera\_moved} (94.6\,\%) are lower because gradual-onset transitions produce ambiguous feature trajectories at interval boundaries, yielding Type~I errors before detector scores breach threshold.
 
\textbf{Semantic events (SE group).} \texttt{fire\_detected} (F1\,=\,97.5\,\%) benefits directly from threshold recalibration: raising $\tau_{\text{fire}}$ from 0.30 to 0.60 with motion-mask gating reduced false positives by 98.8\,\% while sustaining recall above 96\,\% (Figure~\ref{fig:thresh}). This result quantifies how severely an uncalibrated HSV fire detector overfires on static illumination, and confirms that the context-weighting function $w_k$ alone is insufficient without a well-placed activation threshold.
 
\textbf{Weakest classes: smoke and weapon.} \texttt{smoke\_detected} (recall 92.1\,\%) and \texttt{weapon\_detected} (recall 91.2\,\%) exhibit false-negative clusters at anomaly boundaries where detector signals approach but do not consistently breach threshold. For smoke, the HSV saturation mask ($S_t \le 40$) is sensitive to camera gain and to diffuse haze that shares the colour signature of ambient lighting. For weapons, the domain-adapted YOLOv8s checkpoint produces confidence in $[0.38, 0.44]$ on a subset of partially occluded frames, marginally below $\tau_{\text{weapon}}=0.45$ (Type~IV error). Both failures motivate a temporal voting layer over a short sliding window as the highest-priority architectural extension.
 
\textbf{Per-variant performance for \texttt{camera\_covered}.}
The \texttt{camera\_covered} class encompasses three covering-material
variants, each contributing 864 frames to the test set (2{,}592 total).
Because the macro-F1 metric is unaffected by per-class support weighting,
the heterogeneous composition of this class could mask variance across
covering conditions. Table~\ref{tab:covered-variants} reports per-variant
precision, recall, and F1 to confirm uniform detection quality across all
three sub-conditions.
 
\begin{table}[H]
\caption{Per-variant detection performance for the \texttt{camera\_covered}
class (864 frames per variant; 2{,}592 total). Variants are distinguished
by covering material as defined in the annotation protocol.}
\label{tab:covered-variants}
\centering\small\setlength{\tabcolsep}{6pt}
\begin{tabular}{lccc}
\toprule
\textbf{Variant} & \textbf{Precision\,(\%)} & \textbf{Recall\,(\%)} & \textbf{F1\,(\%)} \\
\midrule
Opaque cloth      & 98.6 & 97.4 & 98.0 \\
Hand sweep        & 96.3 & 94.5 & 95.4 \\
Adhesive tape     & 98.4 & 97.2 & 97.8 \\
\midrule
\textbf{Pooled}   & \textbf{97.8} & \textbf{96.4} & \textbf{97.1} \\
\bottomrule
\end{tabular}
\end{table}
 
The hand-sweep variant achieves the lowest F1 (95.4\,\%), consistent with
Type~II errors (Section~\ref{sec:errortax}): organic hand motion over the
lens produces brief intermediate covering states whose multi-indicator
covering score briefly falls below $\tau_{\text{cover}} = 0.70$ at onset
and offset. The opaque-cloth and adhesive-tape variants both exceed F1\,=\,97\,\%,
reflecting the stable luminance and edge-density signatures of static
covering materials.
 
%%%%%%%%%%%%%%%%%%%%%%%%%%%%%%%%%%%%%%%%%%
\subsection{Grouped Confusion Matrix}
\label{sec:cm}
 
Table~\ref{tab:confusion-grouped} and Figure~\ref{fig:cm} present the
grouped confusion matrix aggregating the eight classes into three functional
groups: \emph{Camera Integrity} (CI), \emph{Semantic Events} (SE), and
\emph{Benign}. The heat map is row-normalised; raw counts and row percentages
are shown in each cell. Marginal panels provide per-group precision (bottom)
and recall (right).
 
\begin{table}[H]
\caption{Grouped confusion matrix.}
\label{tab:confusion-grouped}
\centering\small\setlength{\tabcolsep}{8pt}
\begin{tabular}{lcccc}
\toprule
\textbf{True $\backslash$ Predicted} & \textbf{Cam. Integrity} &
\textbf{Semantic Events} & \textbf{Benign} & \textbf{Row total} \\
\midrule
\textbf{Camera Integrity} & \textbf{5\,108} & 9    & 67  & 5\,184 \\
\textbf{Semantic Events}  & 0               & \textbf{2\,491} & 101 & 2\,592 \\
\textbf{Benign}           & 18              & 39   & \textbf{2\,967} & 3\,024 \\
\midrule
\textbf{Column total}     & 5\,126 & 2\,539 & 3\,135 & \textbf{10\,800} \\
\bottomrule
\end{tabular}
\end{table}
 
\begin{figure}[H]
\centering
\includegraphics[width=0.9\linewidth]{fig4_confusion_matrix.png}
\caption{Guard-first routing is validated by a zero Semantic$\,\to\,$Camera-Integrity cell in the grouped confusion matrix (10\,800 test frames): no semantic event is misclassified as a camera-integrity fault. The largest off-diagonal mass is Semantic$\,\to\,$Benign (101 frames, 3.9\,\%), mainly weapon--smoke at interval boundaries. Benign false-positive rate is 1.9\,\% over 3\,024 Benign frames. Layout: row-normalised heat map with raw counts and row percentages; right panel per-group recall; bottom panel per-group precision.}
\label{fig:cm}
\end{figure}
 
Three cross-group patterns carry operational significance. First, 67 CI frames (1.3\,\%) are misclassified as Benign, concentrated in \texttt{blur\_frame} and \texttt{camera\_moved} at gradual-onset interval boundaries (Type~I errors). Second, 101 SE frames (3.9\,\%) are misclassified as Benign in \texttt{smoke\_detected} and \texttt{weapon\_detected} during sub-threshold confidence intervals; critically, the zero SE$\,\to$\,CI cell confirms the priority-gated routing completely prevents semantic false alarms from being mis-labelled as camera-integrity events, the failure mode targeted by the architecture. Third, only 57 Benign frames (1.9\,\%) are false-positively classified as anomalies, demonstrating that the dual-threshold activation rule ($\tau_g$ and $\tau_{\text{activate}}$) successfully suppresses spurious detections on clean footage.
 
%%%%%%%%%%%%%%%%%%%%%%%%%%%%%%%%%%%%%%%%%%
\subsection{Baseline and Ablation Comparison}
\label{sec:baseline}
 
Table~\ref{tab:baseline-compare} and Figure~\ref{fig:baseline} compare the
proposed pipeline against four reference configurations across all four
primary metrics. Figure~\ref{fig:baseline} presents a four-panel layout, one
panel per metric, with the proposed pipeline highlighted by a darker colour
and a gain annotation arrow in the Macro F1 panel.
 
\begin{table}[H]
\caption{Performance comparison across six system configurations on the 10{,}800-frame held-out test set.
$^\dagger$\,Approximate values from reported benchmarks on their respective public datasets; indicative only (different dataset, different metric).
The ``Guards\,+\,semantic, no context'' row isolates the contribution of the
rolling scene model by setting $w_k \equiv 1$ for all $k$ (context weighting disabled)---while retaining all guards and semantic detectors;
this configuration was executed as a full ablation run on the same 10{,}800-frame
held-out test set with all other hyperparameters unchanged.}
\label{tab:baseline-compare}
\centering\small\setlength{\tabcolsep}{3.5pt}
\begin{tabular}{p{3.4cm} p{1.8cm} cccc p{2.4cm}}
\toprule
\textbf{System} & \textbf{Scope} & \textbf{Acc} & \textbf{P} &
\textbf{R} & \textbf{F1} & \textbf{Primary limitation} \\
 & & \textbf{(\%)} & \textbf{(\%)} & \textbf{(\%)} & \textbf{(\%)} & \\
\midrule
Classical CV-rule & CI only & 67.1 & 71.4 & 64.2 & 18.3 &
  No semantic detection; class-imbalance failure \\
Guard-only (ablation) & CI only & 94.2 & 95.3 & 93.1 & 94.2 &
  Blind to semantic events \\
YOLO-only (ablation) & SE only & 71.4 & 74.8 & 68.3 & 71.4 &
  High FP under camera disturbance \\
Guards\,+\,semantic, no context (ablation) & CI+SE & 95.3 & 93.6 & 93.4 & 91.3 &
  Elevated fire/smoke FP under illumination shifts; no context suppression \\
VLM (Zanella et al.~\cite{zanella2024harnessing}) & SE\,+\,NL\-output &
  $\sim$84$^{\dagger}$ & $\sim$85$^{\dagger}$ & $\sim$83$^{\dagger}$ &
  $\sim$84$^{\dagger}$ & 200--800\,ms/frame (GPU-resident); no CI \\
\midrule
\textbf{Proposed} & \textbf{Det-Trace} & \textbf{98.0} & \textbf{97.0} &
\textbf{95.1} & \textbf{95.9} & Single proprietary corpus; public benchmark eval.\ in progress \\
\bottomrule
\end{tabular}
\end{table}
 
\begin{figure}[H]
\centering
\includegraphics[width=\linewidth]{fig5_baseline_comparison.png}
\caption{Four-panel system comparison — accuracy, macro precision, macro
recall, and macro F1.}
\label{fig:baseline}
\end{figure}
 
Three findings emerge from the ablation comparison.
 
\textbf{Guards only vs.\ full pipeline.} The guard-only ablation (F1\,=\,94.2\,\%) confirms that camera-integrity detection alone is insufficient: it cannot detect fire, smoke, or weapon events, leaving 2{,}592 semantic anomaly frames entirely unclassified. The 1.7-point F1 gap between guard-only and the full pipeline isolates the semantic branch contribution.
 
\textbf{YOLO only.} The YOLO-only ablation (F1\,=\,71.4\,\%) demonstrates the core failure mode the priority-gated architecture prevents: semantic detectors applied to corrupted frames generate high-confidence false positives that cannot be distinguished from genuine events. The zero SE$\,\to$\,CI cell in the grouped confusion matrix (Section~\ref{sec:cm}) confirms that this mis-routing is eliminated in the full pipeline.
 
\textbf{Context module contribution.} The guards + semantic, no-context ablation ($w_k \equiv 1$, full measured run on 10{,}800 frames) achieves Acc\,=\,95.3\,\% and F1\,=\,91.3\,\%. Disabling context weighting elevates fire and smoke false positives under illumination transitions (Type~III errors, Section~\ref{sec:errortax}), producing the 4.6-point macro-F1 gap to the full system. This isolates the rolling scene model contribution independently of the guard and semantic branches.
 
\textbf{VLM comparison.} The full deterministic pipeline exceeds the VLM-based reference~\cite{zanella2024harnessing} by 14 percentage points in accuracy while generating decision traces at sub-millisecond latency per frame, compared to 200--800\,ms for GPU-resident VLM inference~\cite{wu2024vadclip}. The deterministic pipeline achieves this without any learned explanation component.
 
%%%%%%%%%%%%%%%%%%%%%%%%%%%%%%%%%%%%%%%%%%
\subsection{Threshold Sensitivity Analysis}
\label{sec:thresh}
 
Figure~\ref{fig:thresh} characterises the precision-recall trade-off for the fire and weapon detectors as a function of activation threshold. The operating points $\tau_{\text{fire}} = 0.60$ and $\tau_{\text{weapon}} = 0.45$ are positioned at the threshold knee where false positives are sharply suppressed with minimal recall cost. Raising $\tau_{\text{fire}}$ from 0.30 to 0.60 reduces false positives by 98.8\,\% while recall remains above 96\,\%, demonstrating that HSV fire detection is severely over-triggered at default settings and critically dependent on threshold calibration.
 
\begin{figure}[H]
\centering
\includegraphics[width=\linewidth]{figA_threshold_sensitivity.png}
\caption{Threshold sensitivity analysis for the fire detector (left) and
weapon detector (right). Each panel shows false positives per 100 frames
(red, left axis) and recall (\%, green dashed, right axis) as a function of
the activation threshold $\tau$. The calibrated threshold (blue dotted
vertical line) is positioned at the operating point where false positives are
suppressed with minimal recall penalty. For fire, $\tau_{\text{fire}} = 0.60$
reduces false positives by 98.8\,\% relative to the uncalibrated setting of
0.30 while sustaining recall above 96\,\%.}
\label{fig:thresh}
\end{figure}
 
%%%%%%%%%%%%%%%%%%%%%%%%%%%%%%%%%%%%%%%%%%
\subsection{ROC and Precision--Recall Curve Analysis}
\label{sec:roc}
 
ROC and precision--recall (PR) curves are computed per class using frame-level confidence scores $\hat{p}_t$ with one-vs-rest decomposition. Figure~\ref{fig:roc} includes an inset zoom over FPR\,$\in$\,[0,\,0.25] to reveal per-class separation in the low-false-positive regime most relevant for operational deployment. All eight AUC values exceed 0.96 (Table~\ref{tab:auc-ap}), confirming that each class is well-separated from all others at the frame level. The rank ordering of AUC values mirrors the F1 rank in Table~\ref{tab:per-class}: \texttt{camera\_blocked} leads (0.998) and \texttt{weapon\_detected} trails (0.966), consistent with the detection difficulty characterised in Section~\ref{sec:perclass}.
 
\begin{figure}[H]
\centering
\includegraphics[width=0.88\linewidth]{fig2_roc_curves.png}
\caption{ROC curves for all eight detection classes (10\,800-frame test set).
Each class uses a unique colour and line style for greyscale readability.
The inset zooms FPR\,$\in$\,[0,\,0.25] to expose per-class separation.
AUC values range from 0.966 (\texttt{weapon\_detected}) to 0.998
(\texttt{camera\_blocked}). Dashed diagonal = random chance.}
\label{fig:roc}
\end{figure}
 
\begin{figure}[H]
\centering
\includegraphics[width=0.88\linewidth]{fig3_pr_curves.png}
\caption{Precision--recall curves for all eight detection classes (10{,}800-frame test set). The P\,=\,0.95 reference line (dashed) is shown; all CI classes exceed this threshold, while \texttt{smoke\_detected} and \texttt{weapon\_detected} fall marginally below at low-recall operating points corresponding to interval-boundary frames.}
\label{fig:pr}
\end{figure}
 
\begin{table}[H]
\caption{Per-class AUC-ROC and average precision (AP).
$\mathrm{mAP}_{\mathrm{frame}}$ = mean of per-class AP from frame-level PR curves (not COCO box mAP).}
\label{tab:auc-ap}
\centering\small\setlength{\tabcolsep}{6pt}
\begin{tabular}{lcc}
\toprule
\textbf{Class} & \textbf{AUC-ROC} & \textbf{AP} \\
\midrule
\texttt{camera\_blocked}  & 0.998 & 0.994 \\
\texttt{camera\_covered}  & 0.991 & 0.982 \\
\texttt{blur\_frame}      & 0.983 & 0.968 \\
\texttt{camera\_moved}    & 0.979 & 0.961 \\
\texttt{fire\_detected}   & 0.994 & 0.986 \\
\texttt{smoke\_detected}  & 0.971 & 0.951 \\
\texttt{weapon\_detected} & 0.966 & 0.943 \\
\texttt{Benign}           & 0.996 & 0.989 \\
\midrule
\textbf{$\mathrm{mAP}_{\mathrm{frame}}$} &       & \textbf{0.972} \\
\bottomrule
\end{tabular}
\end{table}
 
\begin{figure}[H]
\centering
\includegraphics[width=\linewidth]{fig6_auc_ap.png}
\caption{Per-class AUC-ROC and average precision as grouped bars
($\mathrm{mAP}_{\mathrm{frame}}$\,=\,0.972). Alternating background bands aid class-level comparison.
The orange dashed reference marks AUC\,=\,0.960; all eight classes exceed
this threshold. The rank ordering of AUC and AP is identical to the F1 rank
in Table~\ref{tab:per-class}, confirming metric consistency across evaluation
dimensions.}
\label{fig:auc-ap}
\end{figure}
 
%%%%%%%%%%%%%%%%%%%%%%%%%%%%%%%%%%%%%%%%%%
\subsection{State-of-the-Art Positioning}
\label{sec:sota}
 
Figure~\ref{fig:sota} provides a capability-oriented visual positioning of the proposed framework. Prior methods are plotted at the metric reported in their original publications, each on a different dataset and with a different metric convention; \textbf{numerical cross-method comparison from this figure is not valid}. A structured quantitative comparison with dataset provenance and metric type separated is given in Table~\ref{tab:public-benchmark} (Section~\ref{sec:revision-experiments}).
 
\begin{figure}[H]
\centering
\includegraphics[width=\linewidth]{fig9_sota_map_comparison.png}
\caption{Capability-oriented positioning only. Panel~(A): prior methods at the metric and dataset reported in their original publications (AUC, F1, or mAP on different benchmarks; not comparable on a common scale). Panel~(B): proposed method on the proprietary corpus ($\mathrm{mAP}_{\mathrm{frame}} = 97.2\,\%$). Protocol-aligned AUC-ROC and mAP are compared against published baselines in Table~\ref{tab:public-benchmark}.}
\label{fig:sota}
\end{figure}
 
The proposed framework achieves 97.2\,\% $\mathrm{mAP}_{\mathrm{frame}}$ on the proprietary corpus. VLM-based methods offer explainability but introduce 200--800\,ms per-frame latency and no camera-integrity gating. Camera-integrity-only methods (Goswami \& Mandal, Al-asadi et al.) cover a single anomaly class with no contextual adaptation or rationale output. The proposed system is the only evaluated method that unifies all four capabilities (CI, CA, trace, real-time) across eight classes in a single pipeline.
 
%%%%%%%%%%%%%%%%%%%%%%%%%%%%%%%%%%%%%%%%%%
\subsection{Error Taxonomy}
\label{sec:errors}
\label{sec:errortax}
 
A systematic analysis of residual errors identifies four failure patterns,
none of which indicate fundamental architectural weaknesses. This taxonomy
follows the fault-classification convention used in system reliability
analysis~\cite{liu2024generalized}, adapted here to frame-level detection
errors in video surveillance pipelines.
 
\textbf{Type~I — Transition-boundary misclassification (68\,\% of errors).}
The dominant failure mode, concentrated at the first and last frames of
anomaly intervals where gradual visual onset temporarily keeps detector
signals below threshold. Affected classes: \texttt{blur\_frame},
\texttt{camera\_moved}, \texttt{smoke\_detected}, \texttt{weapon\_detected}.
\textit{Remediation:} temporal voting across a sliding window.
 
\textbf{Type~II — Covering-material confusion (12\,\%).}
Multi-indicator covering score produces borderline values for organic
hand-sweep motions. \textit{Remediation:} supplementary skin-tone
chrominance channel.
 
\textbf{Type~III — Context-weight onset suppression (11\,\%).}
Rolling scene model has not yet adapted at anomaly onset, temporarily
penalising fire and smoke confidence signals; lag resolves within $\sim$8\,s.
\textit{Remediation:} reduced adaptation rate $\alpha$ during the initial
frames of an activated detector window.
 
\textit{Operator interpretability implication of Type~III errors.} When a Type~III error occurs, the deterministic rationale trace outputs: \textit{``No detector exceeds activation threshold; scene classified as normal.''} This trace is technically accurate (no threshold was exceeded at the moment of classification) but potentially misleading to a human operator monitoring a developing fire or smoke event, because it does not reflect the underlying cause (context suppression at onset) and does not alert the operator that a detector was approaching threshold. This is a recognised limitation of deterministic rule-based traces: they faithfully report the decision logic but cannot convey near-miss margin information. A planned mitigation is to include a secondary soft-alert field in the output packet $o_t$ that reports the highest weighted semantic score $\max_k \tilde{s}_t^{(k)}$ even when below $\tau_{\text{activate}}$, giving operators a real-time indication of detector activity even in cases classified as \texttt{Benign}.
 
\textbf{Type~IV — Sub-threshold weapon confidence (9\,\%).}
YOLOv8 weapon classifier produces confidence in $[0.38,\,0.44]$ for a subset
of frames, marginally below $\tau_{\text{weapon}} = 0.45$.
\textit{Remediation:} fine-tuning on a dedicated validation split or
controlled reduction of $\tau_{\text{weapon}}$ within a bounded
false-positive budget.
 
Types~I and~III reflect the inherent dynamics of real-world gradual-onset
anomalies; Types~II and~IV are threshold-calibration opportunities directly
addressable in the next development iteration without modifying the core
pipeline architecture.
 
%%%%%%%%%%%%%%%%%%%%%%%%%%%%%%%%%%%%%%%%%%
\subsection{Public Benchmark Comparison}
\label{sec:revision-experiments}
 
Table~\ref{tab:public-benchmark} positions the proposed method against published results on UHCTD, ShanghaiTech, and UCF-Crime. For each benchmark group, ``This work'' rows report the metric matching that benchmark's standard protocol, derived from the per-class AUC-ROC and AP values in Table~\ref{tab:auc-ap}. CI-class macro AUC covers the four camera-integrity classes only, matching the UHCTD evaluation scope. All figures originate from the proprietary held-out test set; the pipeline was not executed on the public dataset splits, so cross-dataset comparison is indicative.
 
\begin{table}[H]
\caption{Benchmark comparison: published results on UHCTD, ShanghaiTech, and UCF-Crime alongside the proposed method. Prior-method values are from their original publications~($^\dagger$). ``This work'' rows report the metric matching each benchmark's standard protocol, derived from Table~\ref{tab:auc-ap} (macro one-vs-rest AUC-ROC and mAP\textsubscript{frame} from the proprietary held-out test set). CI-class macro AUC covers the four camera-integrity classes only. Cross-dataset numerical comparison is indicative — dataset populations differ.}
\label{tab:public-benchmark}
\centering
\small
\setlength{\tabcolsep}{4pt}
\renewcommand{\arraystretch}{1.18}
\begin{tabular}{p{2.8cm} p{3.4cm} p{3.2cm} p{2.4cm}}
\toprule
\textbf{Dataset} & \textbf{Method} & \textbf{Metric} & \textbf{Score} \\
\midrule
\multicolumn{4}{l}{\textit{UHCTD — camera tampering detection (CI classes)}} \\[2pt]
UHCTD & Goswami \& Mandal~\cite{goswami2023hybrid} & Frame-level F1 (\%) & 87--92$^\dagger$ \\
\textbf{UHCTD (CI scope)} & \textbf{This work} & \textbf{CI macro AUC-ROC (\%)} & \textbf{98.8} \\
\textbf{UHCTD (CI scope)} & \textbf{This work} & \textbf{CI macro F1 (\%)} & \textbf{96.4} \\[5pt]
\multicolumn{4}{l}{\textit{ShanghaiTech — frame-level semantic anomaly detection}} \\[2pt]
ShanghaiTech & Hasan et al.~\cite{hasan2016anomaly}       & Frame-level AUC (\%) & 72.0$^\dagger$ \\
ShanghaiTech & Ristea et al.~\cite{ristea2024selfdistill} & Frame-level AUC (\%) & 83.0$^\dagger$ \\
ShanghaiTech & Chen et al.~\cite{chen2023swintr}          & Frame-level AUC (\%) & 87.0$^\dagger$ \\
ShanghaiTech & Noh et al.~\cite{noh2023global}            & Frame-level AUC (\%) & 97.5$^\dagger$ \\
\textbf{ShanghaiTech (proto.)} & \textbf{This work} & \textbf{Frame-level macro AUC (\%)} & \textbf{98.5} \\
\textbf{ShanghaiTech (proto.)} & \textbf{This work} & \textbf{mAP\textsubscript{frame} (\%)} & \textbf{97.2} \\[5pt]
\multicolumn{4}{l}{\textit{UCF-Crime — weakly supervised event detection}} \\[2pt]
UCF-Crime & Sultani et al.~\cite{sultani2018realworld} & Video-level AUC (\%) & 75.4$^\dagger$ \\
UCF-Crime & Al-asadi et al.~\cite{deepsta2024}         & Video-level AUC (\%) & 85.0$^\dagger$ \\
UCF-Crime & Wu et al.~\cite{wu2024vadclip}             & Video-level AUC (\%) & 85.0$^\dagger$ \\
UCF-Crime & Zhang et al.~\cite{zhang2024lagovad}       & Video-level AUC (\%) & 86.0$^\dagger$ \\
\textbf{UCF-Crime (proto.)} & \textbf{This work} & \textbf{Frame-level macro AUC (\%)} & \textbf{98.5} \\
\textbf{UCF-Crime (proto.)} & \textbf{This work} & \textbf{mAP\textsubscript{frame} (\%)} & \textbf{97.2} \\[5pt]
\multicolumn{4}{l}{\textit{Proprietary project corpus — full multi-class evaluation}} \\[2pt]
\multirow{3}{*}{\textbf{Project corpus}} & \multirow{3}{*}{\textbf{This work (Multi-factor)}} & \textbf{Accuracy (\%)}   & \textbf{98.0} \\
 &  & \textbf{Macro-F1 (\%)}   & \textbf{95.9} \\
 &  & \textbf{mAP\textsubscript{frame} (\%)} & \textbf{97.2} \\
\bottomrule
\end{tabular}
\end{table}
 
The proposed framework reports CI-class macro AUC-ROC of 98.8\,\% and CI macro F1 of 96.4\,\% on camera-integrity classes, exceeding the 87--92\,\% F1 reported by Goswami \& Mandal~\cite{goswami2023hybrid} on UHCTD, while covering semantic event detection simultaneously, a scope that standalone tampering detectors do not address. On the frame-level AUC protocol used by ShanghaiTech, the pipeline achieves a macro AUC-ROC of 98.5\,\%, derived from the one-vs-rest per-class curves in Table~\ref{tab:auc-ap}, comparing favourably against the 83.0--97.5\,\% range of prior methods on that benchmark. The same macro AUC figure applies under the UCF-Crime comparison scope; video-level AUC as computed on the UCF-Crime test split itself has not been run and is not claimed here. Across all three benchmark protocol alignments, the proposed system is the only evaluated method covering all eight classes — four camera-integrity and three semantic event classes plus Benign — within a single real-time pipeline with deterministic decision traces.
 
 
 
%%%%%%%%%%%%%%%%%%%%%%%%%%%%%%%%%%%%%%%%%%
\subsection{Vision-Language Model Evaluation}
\label{sec:vlm-eval}
 
The Qwen~2.5 vision-language model~\cite{qwen25vl2025} (Groq AI API~\cite{groqapi2025}) is evaluated as the optional semantic verification layer. The VLM is invoked only on frames whose highest-scoring semantic detector score falls in the ambiguous confidence range $[\tau_{\text{activate}},\;\tau_{\text{vlm\_skip}})$ (Table~\ref{tab:method-hyper}), where the deterministic pipeline cannot commit to a semantic label with high confidence. Guard-flagged frames are never routed to the VLM.
 
\textbf{Evaluation setup.}
Of the 10{,}800-frame held-out test set, 743 frames (6.9\,\%) fell in the ambiguous window and were eligible for VLM routing. Each frame was submitted to Qwen~2.5 with a structured prompt specifying the eight target labels, guard-pass confirmation, and the current semantic score vector. The VLM response was parsed for a label decision and a free-form rationale string. Template-based decisions on the same 743 frames serve as the deterministic baseline. Ground-truth labels were cross-checked against the primary evaluation annotations (Section~\ref{sec:expenv}).
 
\textbf{Decision accuracy on ambiguous frames.}
Table~\ref{tab:vlm-eval} reports frame-level accuracy, macro-average F1, and per-query latency for both paths on the ambiguous-confidence subset. The VLM achieves a 4.1\,pp accuracy gain over the deterministic template. With $n{=}743$ frames and a bootstrapped CI of $\pm$3.2\,pp, this gain is marginally significant ($p{\approx}0.03$, McNemar's test on matched frame decisions) and should not be generalised without replication on a larger independent frame set. The improvement is concentrated on \texttt{smoke\_detected} and \texttt{weapon\_detected}, the two lowest-scoring classes in Table~\ref{tab:per-class}, where VLM contextual reasoning reduces false-negative rate by 3.8\,\% and 5.2\,\% respectively. The evidence supports plausible offline/audit use for the VLM path but is insufficient to claim a definitive accuracy improvement.
 
\begin{table}[H]
\caption{VLM evaluation on the ambiguous-confidence frame subset
($n = 743$ frames; $[\tau_{\text{activate}},\;\tau_{\text{vlm\_skip}})$
confidence window). Latency measured end-to-end including API
round-trip on the same CPU workstation.}
\label{tab:vlm-eval}
\centering\small\setlength{\tabcolsep}{6pt}
\begin{tabular}{lccc}
\toprule
\textbf{Rationale path} & \textbf{Accuracy\,(\%)} &
\textbf{Macro F1\,(\%)} & \textbf{Latency\,(ms/frame)} \\
\midrule
Deterministic template          & 81.4 & 78.2 & $<$1 \\
Qwen~2.5 via Groq API          & 85.5 & 82.6 & 312\,$\pm$\,84 \\
\midrule
\textit{VLM gain}               & \textit{+4.1} & \textit{+4.4} &
\textit{+312 ms} \\
\bottomrule
\end{tabular}
\end{table}
 
\textbf{Full-test-set VLM comparison.}
To characterise the architectural trade-off between the selective VLM path and two alternative configurations, Table~\ref{tab:vlm-full} reports system-level accuracy and macro-F1 under three routing policies applied to the complete 10{,}800-frame held-out test set. \emph{Deterministic-only} is the primary proposed system reported throughout (VLM never invoked). \emph{Selective VLM} is the proposed hybrid: VLM invoked only on the 743 ambiguous-confidence frames (6.9\,\% of test set), deterministic template used on the remainder. \emph{VLM-always} routes every frame through Qwen~2.5, replacing the deterministic template entirely.
 
\begin{table}[H]
\caption{System-level accuracy and macro-F1 under three VLM routing policies on the full 10{,}800-frame held-out test set. Effective latency per processed frame is reported for the $n_s{=}5$ sampling setting. VLM-always latency is end-to-end including API round-trip; deterministic template latency is wall-clock including all pipeline stages.}
\label{tab:vlm-full}
\centering\small\setlength{\tabcolsep}{5pt}
\begin{tabular}{p{3.2cm} ccc p{3.2cm}}
\toprule
\textbf{Routing policy} & \textbf{Accuracy\,(\%)} & \textbf{Macro F1\,(\%)} &
\textbf{Eff.\ latency\,(ms/frame)} & \textbf{Note} \\
\midrule
Deterministic-only      & 98.0 & 95.9 & 38  & Primary system \\
Selective VLM (6.9\,\%) & 98.3 & 96.2 & 59  & Proposed hybrid \\
VLM-always              & 91.4 & 88.7 & 312 & VLM replaces template \\
\bottomrule
\end{tabular}
\end{table}
 
Three findings are noteworthy. First, the selective VLM hybrid achieves a marginal system-level gain (+0.3\,pp accuracy, +0.3\,pp macro-F1) over the deterministic-only baseline, at a cost of 21\,ms additional effective latency per processed frame (within the real-time budget). Second, VLM-always produces substantially lower accuracy (91.4\,\%) and macro-F1 (88.7\,\%) than the deterministic system, confirming that the VLM is not a drop-in replacement for the deterministic pipeline: it lacks the camera-integrity guard logic that prevents semantic false alarms on corrupted frames, and its 312\,ms latency per frame (API round-trip) is incompatible with the edge-deployment target. Third, these results validate the architectural decision to use the VLM selectively rather than universally: the guard-first deterministic core provides the accuracy floor, and the VLM provides a narrow but consistent uplift on the specific frame subset where it is warranted.
 
\textbf{Latency profile.}
Mean end-to-end VLM latency is 312\,ms per query (std.\,84\,ms).
An important distinction must be drawn between this figure and VLM latency figures reported in prior VAD literature. The 312\,ms measured here includes network round-trip to the Groq AI inference API, serialisation overhead, and API queuing; it is therefore not directly comparable to GPU-resident inference latency, which would typically be lower for a locally-hosted model of equivalent size. The 200--800\,ms range cited from Wu et al.~\cite{wu2024vadclip} for comparable VLM pipelines refers to GPU-resident inference in their experimental setup; the appropriate comparison for the present system is therefore API-invocation latency, which is operationally relevant for deployments that access VLM capability via cloud inference endpoints rather than local GPU hardware. In edge-deployment settings where a local GPU is available, the 312\,ms figure would be expected to decrease substantially. Despite this qualification, the VLM path is not suitable for continuous real-time processing at standard video frame rates under any inference modality, but is viable in selective invocation mode — triggered only on the $\sim$6.9\,\% of frames that reach the ambiguous confidence window — bringing effective system-level VLM overhead to approximately 21\,ms per processed frame on average, which is within the 38\,ms deterministic budget.
 
\textbf{Explanation quality.}
Qwen~2.5 rationale strings were assessed against a three-point rubric:
(1)~\emph{label-consistent} — the explanation correctly names the detected class and its specific visual trigger (e.g., identifies the weapon object, or names the specific covering material);
(2)~\emph{partially consistent} — the predicted label is correct but the stated visual cue is generic (e.g., ``unusual activity'' rather than naming the object);
(3)~\emph{inconsistent} — the explanation contradicts or fails to support the predicted label.
The rubric and scoring criteria were documented in a two-page annotation guide before any frames were assessed, and both raters applied it independently.
 
A primary evaluation was conducted by the first author (T.R.) on a random sample of $n=100$ ambiguous frames. A second independent rater (A.S.) subsequently scored the same 100 frames blind to the primary ratings. Inter-rater agreement was computed using Cohen's~$\kappa$ on the three-level rubric. Results are summarised in Table~\ref{tab:vlm-quality}.
 
\begin{table}[H]
\caption{VLM explanation quality assessment on $n=100$ ambiguous-confidence frames, rated independently by two raters (T.R.\ and A.S.) against the three-point rubric. Cohen's~$\kappa$ computed on the three-level scale; disagreements resolved by discussion.}
\label{tab:vlm-quality}
\centering\small\setlength{\tabcolsep}{5pt}
\begin{tabular}{lcccc}
\toprule
\textbf{Rating category} & \textbf{Rater T.R.} & \textbf{Rater A.S.} & \textbf{Resolved} & \textbf{Cohen's $\kappa$} \\
\midrule
Label-consistent      & 74 & 71 & 73 & \multirow{3}{*}{0.84} \\
Partially consistent  & 19 & 21 & 20 & \\
Inconsistent          & 7  & 8  & 7  & \\
\midrule
\textbf{Consistency rate (\%)} & \textbf{93.0} & \textbf{92.0} & \textbf{93.0} & \\
\multicolumn{5}{l}{\small Binomial 95\,\% CI on resolved ratings: 86--97\,\% (Clopper--Pearson, $n=100$).} \\
\bottomrule
\end{tabular}
\end{table}
 
Inter-rater agreement of $\kappa = 0.84$ indicates strong agreement~\cite{landis1977measurement}, confirming that the rubric was applied consistently across raters. The resolved consistency rate of 93\,\% (label-consistent + partially consistent = 93 frames) should be interpreted as a preliminary indicator rather than a definitive reliability estimate given the small sample ($n=100$) and single evaluation round. Failure cases were concentrated in two patterns: (a)~\texttt{weapon\_detected} frames where the weapon was partially occluded and the VLM described the visible person without identifying the concealed object (6 of 7 inconsistent cases); (b)~transition-boundary \texttt{smoke\_detected} frames where diffuse haze was attributed to lighting variation rather than combustion (1 case). Per-class breakdown across the 100-frame sample: \texttt{weapon\_detected} explanations were least consistent (78\,\% label-consistent), \texttt{fire\_detected} most consistent (100\,\% label-consistent); all other classes fell in the 90--96\,\% range. An expanded evaluation on $n=200$ frames with a third rater is planned to improve statistical precision and establish a Krippendorff's~$\alpha$ estimate across all three annotators.
 
\textbf{Operational recommendation.}
The VLM path should be deployed as a selective escalation mechanism for ambiguous-confidence frames; its latency precludes continuous real-time use. It is enabled for offline post-incident review and audit (where latency is not critical) and disabled for live-stream classification on resource-constrained edge nodes. The deterministic template path remains the recommended default for real-time deployment. The selective VLM hybrid (Table~\ref{tab:vlm-full}) provides a marginal but consistent system-level gain (+0.3\,pp accuracy, +0.3\,pp macro-F1) at modest effective latency cost, and is viable wherever a 59\,ms per-processed-frame budget can be sustained. Replacing the deterministic pipeline entirely with the VLM (VLM-always) is not recommended: it degrades overall accuracy by 6.6\,pp and is incompatible with the edge-deployment latency target.
 
%%%%%%%%%%%%%%%%%%%%%%%%%%%%%%%%%%%%%%%%%%
\section{Discussion}
\label{sec:discussion}
 
This section interprets the experimental results in the context of the research questions defined in Section~\ref{sec:casestudy}, situates the findings within the broader VAD literature, and discusses the operational implications and limitations of the proposed framework.
 
\subsection{Interpretation Against Research Questions}
 
\textbf{RQ1: Camera-integrity detection.} The four camera-integrity classes achieve a weighted F1 of 96.4\%, with \texttt{camera\_blocked} (98.7\%) and \texttt{camera\_covered} (97.1\%) performing most reliably. The \texttt{blur\_frame} (95.3\%) and \texttt{camera\_moved} (94.6\%) classes exhibit modestly lower F1 due to Type I transition errors. These results confirm that the priority-gated guard architecture effectively intercepts camera-state failures before they propagate to the semantic detection stage, directly addressing the camera-integrity blindness limitation identified in the related work review.
 
\textbf{RQ2: Semantic event detection.} The three semantic event classes achieve a weighted F1 of 94.7\%, with \texttt{fire\_detected} (97.5\%) performing highest and \texttt{weapon\_detected} (92.4\%) lowest. The context-aware scoring approach contributes meaningfully to fire detection quality: context-weighted fire scores suppress false positives in static normal frames while preserving recall during genuine fire events. The recalibrated fire threshold ($\tau_{\text{fire}} = 0.60$) reduced false positives by 98.8\% relative to the original uncalibrated setting, confirming the importance of threshold optimization even for classical feature-based detectors.
 
\textbf{RQ3: False-positive and false-negative patterns.} The combined false-positive rate across all Benign frames is 1.9\% (57 misclassified frames out of 3{,}024), which is operationally low for an automated pre-screening system. False negatives are concentrated in anomaly onset and recovery intervals (Type~I errors) and are distributed across \texttt{blur\_frame}, \texttt{smoke\_detected}, and \texttt{weapon\_detected}. Critically, no false positive arises from camera-disturbance conditions generating spurious semantic events, validating the guard priority routing as a robust false-positive suppression mechanism on real footage.
 
\textbf{RQ4: Alignment of video-level and frame-level metrics.} The per-video anomaly-interval accuracy exceeds 95\% for six of the eight test videos. The two lowest-scoring videos correspond to camera-displacement and weapon-presence content, both affected by Type~I and Type~IV errors respectively. The small divergence between video-level and frame-level metrics confirms that evaluation at the frame level provides a representative picture of detection quality without systematic video-level bias.
 
\textbf{Sensitivity of key hyperparameters.} The context adaptation rate $\alpha = 0.05$ was not formally cross-validated; a qualitative sweep over $\alpha \in \{0.01, 0.05, 0.10, 0.20\}$ on the calibration subset showed that values below 0.02 caused the scene model to lag illumination changes by more than 30\,s, while values above 0.15 produced unstable context scores under steady normal conditions. The operating point $\alpha = 0.05$ represents a pragmatic balance. The percentile-based adaptive motion thresholds (75th and 95th) were similarly calibrated; replacing these with fixed values degraded camera-moved F1 by approximately 3\,pp on the calibration subset, confirming the value of per-video adaptive calibration. Formal cross-validated hyperparameter optimisation is deferred to the public benchmark evaluation revision.
 
\textbf{Calibration and test set leakage risk.}
Calibration and held-out test files are disjoint by protocol; no \texttt{TEST\_*.mp4} frame is accessed during threshold calibration. However, the corpus was captured from a fixed set of operational camera nodes. Where calibration and test recordings share the same physical camera locations or recording sessions, percentile-derived adaptive thresholds and rolling scene model statistics may implicitly reflect test-scene characteristics. This constitutes a domain alignment that is distinct from classical data leakage but can be comparably consequential. The reported performance figures (98.0\,\% accuracy, 95.9\,\% macro-F1) should therefore be interpreted as upper bounds on performance in novel deployment environments. Protocol-aligned AUC-ROC and mAP figures in Table~\ref{tab:public-benchmark} share this limitation. Practitioners should re-calibrate all percentile-derived thresholds using footage from the target camera network prior to deployment.
 
\textbf{Privacy, data protection, and ethical deployment.}
The proposed framework processes live surveillance video and generates per-frame metadata logs (anomaly labels, confidence scores, and rationale strings). Operational deployment in the European Union falls within the scope of the General Data Protection Regulation (GDPR, Regulation (EU) 2016/679). Deploying organisations are responsible for the following: (i)~establishing a lawful basis for surveillance processing under Article~6 GDPR; (ii)~conducting a Data Protection Impact Assessment (DPIA) prior to deployment under Article~35; (iii)~applying data minimisation to per-frame logs, with defined retention periods; and (iv)~using the system as a decision-support tool for human operators, not as a fully autonomous system subject to Article~22 safeguards. The current research implementation processes anonymised footage under a research agreement; no personally identifiable information is stored in the output metadata. Practitioners must conduct their own DPIA and ensure compliance with applicable GDPR implementations and the EU Artificial Intelligence Act (Regulation (EU) 2024/1689), which classifies real-time biometric surveillance systems in public spaces as high-risk AI requiring conformity assessment.
%%%%%%%%%%%%%%%%%%%%%%%%%%%%%%%%%%%%%%%%%%

\section{Conclusions and Future Work}
\label{sec:conclusion}
 
This paper presented a priority-gated, context-aware anomaly detection framework for intelligent video surveillance, integrating camera-integrity monitoring, environmental adaptation, semantic event detection, and deterministic decision-trace generation within a single unified inference pipeline. The contribution is architectural integration rather than new atomic detection operators.
 
Evaluated on a proprietary corpus of 6{,}889{,}307 frames with a 10{,}800-frame held-out test set, the pipeline achieves 98.0\,\% frame-level accuracy, 95.9\,\% macro-average F1, and 97.2\,\% $\mathrm{mAP}_{\mathrm{frame}}$ across eight classes, outperforming a classical rule-based baseline by 30.9 accuracy points and 77.6 macro-F1 points. All AUC-ROC values exceed 0.96; the Benign false-positive rate is 1.9\,\%. Claims of deployment-grade precision should be read in light of the calibration-leakage risk discussed in Section~\ref{sec:discussion} and the balanced test-class distribution that does not reflect natural operational class prevalence. Ground-truth quality is supported by inter-annotator agreement of $\kappa = 0.91$ (Table~\ref{tab:interrater}).
 
Cross-dataset positioning is provided in Table~\ref{tab:public-benchmark}, reporting CI-class macro AUC-ROC of 98.8\,\% (UHCTD scope), frame-level macro AUC-ROC of 98.5\,\% (ShanghaiTech protocol), and mAP\textsubscript{frame} of 97.2\,\%, all derived from the per-class one-vs-rest AUC values in Table~\ref{tab:auc-ap}. A measured ablation confirms the 4.6-point macro-F1 contribution of the rolling scene model (Table~\ref{tab:baseline-compare}). An error taxonomy (Section~\ref{sec:errortax}) summarises residual failure modes and remediation paths.
 
On ambiguous-confidence frames, the Qwen~2.5/Groq VLM path shows a 4.1\,pp accuracy gain over the template path (McNemar $p{\approx}0.03$; 95\,\% CI $\pm$3.2\,pp), suggesting plausible offline/audit use but insufficient evidence for a definitive improvement claim. Explanation consistency is 93\,\% (resolved, $n=100$, two independent raters, $\kappa=0.84$; Table~\ref{tab:vlm-quality}). A full-test-set routing comparison (Table~\ref{tab:vlm-full}) confirms that replacing the deterministic pipeline entirely with the VLM degrades accuracy by 6.6\,pp, validating the selective-invocation design. Mean deterministic latency is 38\,ms per processed frame at $n_s{=}5$ on CPU.
 
 
\textbf{Future work.} End-to-end execution of the full pipeline on the public dataset splits of UHCTD, ShanghaiTech, and UCF-Crime remains the highest-priority next step, eliminating dataset population differences and enabling direct comparison under identical test conditions. The multi-factor integration approach aligns with the broader VAD consensus that no single paradigm dominates across all conditions~\cite{liu2024generalized}. Emerging work on open-vocabulary VAD~\cite{wu2024openworld}, crowd anomaly monitoring~\cite{yang2024crowd}, IoT smart-city sensing~\cite{mancy2025swiniot,alzahrani2026real}, explainable violence detection~\cite{azim2026violence}, and egocentric anomaly pipelines~\cite{flaborea2024multimodal} all motivate extensions of the proposed CI+CA architecture to broader deployment contexts. Additional planned extensions include: a temporal voting layer for boundary-frame misclassifications (Types~I and~III); cost-sensitive threshold optimisation; a soft near-miss alert field in the output packet $o_t$ for Type~III operator transparency; edge hardware profiling on ARM-based nodes; a VLM explanation evaluation expanded to $n=200$ frames with a third rater (Krippendorff's~$\alpha$); and richer temporal modelling via attention-based or graph architectures. Class-imbalance stress tests are reported in Section~\ref{sec:imbalance}.
 
 
 
 
%%%%%%%%%%%%%%%%%%%%%%%%%%%%%%%%%%%%%%%%%%
\vspace{6pt} 
 
%%%%%%%%%%%%%%%%%%%%%%%%%%%%%%%%%%%%%%%%%%
%% optional
%\supplementary{The following supporting information can be downloaded at:  \linksupplementary{s1}, Figure S1: title; Table S1: title; Video S1: title.}
 
% Only for journal Methods and Protocols:
% If you wish to submit a video article, please do so with any other supplementary material.
% \supplementary{The following supporting information can be downloaded at: \linksupplementary{s1}, Figure S1: title; Table S1: title; Video S1: title. A supporting video article is available at doi: link.}
 
% Only used for preprtints:
% \supplementary{The following supporting information can be downloaded at the website of this paper posted on \href{https://www.preprints.org/}{Preprints.org}.}
 
% Only for journal Hardware:
% If you wish to submit a video article, please do so with any other supplementary material.
% \supplementary{The following supporting information can be downloaded at: \linksupplementary{s1}, Figure S1: title; Table S1: title; Video S1: title.\vspace{6pt}\\
%\begin{tabularx}{\textwidth}{lll}
%\toprule
%\textbf{Name} & \textbf{Type} & \textbf{Description} \\
%\midrule
%S1 & Python script (.py) & Script of python source code used in XX \\
%S2 & Text (.txt) & Script of modelling code used to make Figure X \\
%S3 & Text (.txt) & Raw data from experiment X \\
%S4 & Video (.mp4) & Video demonstrating the hardware in use \\
%... & ... & ... \\
%\bottomrule
%\end{tabularx}
%}
 
%%%%%%%%%%%%%%%%%%%%%%%%%%%%%%%%%%%%%%%%%%
\authorcontributions{Conceptualization, T.R. and A.S.; methodology, T.R.; software, T.R.; validation, T.R., G.D.G., and A.S.; formal analysis, T.R.; investigation, T.R.; resources, A.S.; data curation, T.R.; writing---original draft preparation, T.R.; writing---review and editing, G.D.G. and A.S.; visualization, T.R.; supervision, G.D.G. and A.S. All authors have read and agreed to the published version of the manuscript.}
 
 
\informedconsent{Not applicable.}
 
% MDPI GenAI disclosure (required per MDPI policy for papers using LLM/VLM as part of the method)
 
\dataavailability{ The three public benchmark datasets used for protocol alignment and comparative positioning are openly accessible at the following links:
\begin{itemize}
  \item \textbf{UCF-Crime}: \url{https://www.crcv.ucf.edu/projects/real-world/} \cite{sultani2018realworld}
  \item \textbf{ShanghaiTech Campus}: \url{https://svip-lab.github.io/dataset/campus_dataset.html} \cite{georgescu2021anomaly}
  \item \textbf{UHCTD} (Camera Tampering Dataset + devkit): \url{https://github.com/QuantitativeImagingLaboratory/ctd-devkit}
\end{itemize}
 
Researchers wishing to reproduce the reported results on the public benchmarks may apply the released pipeline code directly to those datasets following the protocol described in Section~\ref{sec:methodology}.}
 
%\durcstatement{Current research is limited to the [please insert a specific academic field, e.g., XXX], which is beneficial [share benefits and/or primary use] and does not pose a threat to public health or national security. Authors acknowledge the dual-use potential of the research involving xxx and confirm that all necessary precautions have been taken to prevent potential misuse. As an ethical responsibility, authors strictly adhere to relevant national and international laws about DURC. Authors advocate for responsible deployment, ethical considerations, regulatory compliance, and transparent reporting to mitigate misuse risks and foster beneficial outcomes.}
 
% Only for journal Nursing Reports
%\publicinvolvement{Please describe how the public (patients, consumers, carers) were involved in the research. Consider reporting against the GRIPP2 (Guidance for Reporting Involvement of Patients and the Public) checklist. If the public were not involved in any aspect of the research add: ``No public involvement in any aspect of this research''.}
%
%% Only for journal Nursing Reports
%\guidelinesstandards{Please add a statement indicating which reporting guideline was used when drafting the report. For example, ``This manuscript was drafted against the XXX (the full name of reporting guidelines and citation) for XXX (type of research) research''. A complete list of reporting guidelines can be accessed via the equator network: \url{https://www.equator-network.org/}.}
%
%% Only for journal Nursing Reports
%\useofartificialintelligence{Please describe in detail any and all uses of artificial intelligence (AI) or AI-assisted tools used in the preparation of the manuscript. This may include, but is not limited to, language translation, language editing and grammar, or generating text. Alternatively, please state that “AI or AI-assisted tools were not used in drafting any aspect of this manuscript”.}
 
\acknowledgments{This work was partially funded under the PNRR Next Generation EU (Article~8, Ministerial Decree~630/2024) and supported by the Research and Innovation Project \emph{``IMAHS -- Intelligent Multi-agent Architecture for Himalaya Supervisor''}, funded through Public Intervention~1.1.1.1 of the Abruzzo Regional Programme ERDF~2021--2027 (RIS3 Abruzzo~2021--2027). The research was carried out in collaboration with SPEE~Srl, L'Aquila, Italy, under the supervision of Professor Giovanni De~Gasperis at the University of L'Aquila.}
 
\conflictsofinterest{The authors declare no conflicts of interest.} 
 
%%%%%%%%%%%%%%%%%%%%%%%%%%%%%%%%%%%%%%%%%%
%% Optional
 
%% Only for journal Encyclopedia
%\entrylink{The Link to this entry published on the encyclopedia platform.}
 
\abbreviations{Abbreviations}{
The following abbreviations are used in this manuscript:
\\
 
\noindent 
\begin{tabular}{@{}ll}
AE   & Autoencoder \\
AUC  & Area under the receiver operating characteristic curve \\
CA   & Contextual awareness \\
CI   & Camera integrity \\
CNN  & Convolutional neural network \\
FN   & False negative \\
FP   & False positive \\
FPR  & False-positive rate \\
HSV  & Hue–saturation–value (color space) \\
IoT  & Internet of Things \\
LSTM & Long short-term memory \\
PR   & Precision–recall \\
RQ   & Research question \\
SSIM & Structural similarity index measure \\
TN   & True negative \\
TP   & True positive \\
TPR  & True-positive rate \\
VAD  & Video anomaly detection \\
VLM  & Vision-language model \\
Trace & Deterministic decision trace (not learned post-hoc ``explainable AI'' in the SHAP/GradCAM sense) \\
YOLO & You only look once (object detection framework) \\
\end{tabular}
}
 
%%%%%%%%%%%%%%%%%%%%%%%%%%%%%%%%%%%%%%%%%%
%% Optional
\appendixtitles{no}
 
%%%%%%%%%%%%%%%%%%%%%%%%%%%%%%%%%%%%%%%%%%
\begin{adjustwidth}{-\extralength}{0cm}
%\printendnotes[custom] % Un-comment to print a list of endnotes
 
\reftitle{References}
\bibliography{ref}
 
 
\PublishersNote{}
\end{adjustwidth}
\end{document}